\documentclass{article}
\usepackage[margin=0.8in]{geometry}
\usepackage{amsmath,amssymb,amsthm,mathtools}
\usepackage{mathrsfs,bm,bbm}
\usepackage{graphicx,booktabs,array,multirow,tabularx,longtable}
\usepackage{enumitem}
\usepackage{xcolor}
\usepackage{algorithm,algpseudocode}
\usepackage{subcaption,tikz}
\usepackage{siunitx,cancel,accents}
\usepackage{microtype}
\usepackage[numbers,sort&compress]{natbib}
\usepackage[colorlinks=true,linkcolor=blue,citecolor=blue,urlcolor=blue]{hyperref}
\usepackage{cleveref}
\setlength{\emergencystretch}{2em}
\newtheorem{assumption}{Assumption}
\newtheorem{definition}{Definition}
\newtheorem{theorem}{Theorem}
\newtheorem{lemma}{Lemma}
\newtheorem{proposition}{Proposition}
\newtheorem{openendedquestion}{Open-ended Question}
\newtheorem{conjecture}{Conjecture}
\newcommand{\coreref}[1]{\ref{#1}}

\begin{document}

\sloppy

\section*{panel\_\allowbreak{}bivariate\_\allowbreak{}cic\_\allowbreak{}critical\_\allowbreak{}transport\_\allowbreak{}inference}

\noindent\textit{Specialization:} v1\par

\paragraph{TL;DR.} The note proves intrinsic flat-torus changes-in-changes identification, the complete translation-aware nested derivative, and the critical statistical theory for the specified positive-normalized periodized-KDE plug-in.  Uniform transformed-kernel Green estimates and the mixed two-endpoint remainder yield the full four-cell covariance, the fixed-grid Gaussian limit at rate \(\sqrt{n/L_n}\), consistent adjoint covariance, simultaneous multiplier ellipses, and the matched \(L_n/n\) squared-risk rate.

\section{Project justification}

\paragraph{Gap.} Published fixed-source transport linearizations do not jointly handle two varying endpoints, a generated target law, torus harmonic translations, and bandwidth-scale Green poles.  The note closes that combined analytic gap on its smooth positive uniformly elliptic selected-no-cut flat-torus class; it does not assert an inclusion relation with the published Euclidean classes.

\paragraph{Niche.} A uniformly coercive weighted-Hodge Schur complement determines both selected translations.  After pulling both transport equations into treated-post coordinates, the only logarithmically singular component is one outer Green response to four independent cell noises, while the second elliptic response and the complete harmonic correction have bounded score norm.

\paragraph{Fill.} The note proves the variable-coefficient frozen-pole replacement, all generated-target within-cell cross terms, opposite-fold adjoint stability, and the translation-aware \(C U_1U_0\) nonlinear bound throughout \(1/[2(s-2)]\leq a_h<1/6\).  These results certify the covariance, Gaussian, ellipse, and upper-risk arguments, while the explicit local alternatives certify the matching Hellinger-local rate lower bound.

\section{Related work}

TorousGunsiliusRigollet2024OTCiC provides the Euclidean proper-convex-support changes-in-changes identification result; the present intrinsic full-support flat-torus theorem is a distinct analogue, not a proved extension or restriction of that class.  GonzalezSanzSheng2024Linearization supplies a fixed-source, one-target density-curve linearization; the present two-endpoint generated-target and harmonic construction is informed by it but mathematically incomparable in domain and input structure.  ManoleBalakrishnanNilesWeedWasserman2023CLT proves the dimension-two critical phenomenon in the uniform fixed-source case.  Here that phenomenon is generalized only on the stronger smooth positive uniformly elliptic selected-no-cut fixed-grid independent-four-sample torus class, and the known-cell uniform reduction recovers \(I_2/(4\pi)\).  The comparisons do not claim nonperiodic-domain coverage or formal class inclusion.

\section{Setup}

Target (estimand / structural parameter): \(\tau(P)=(-a_T-\nabla\phi_T(\widetilde y_j))_{j=1}^J\).
Primary functional / estimator / diagnostic: Candidate critical inference target: \(\widehat\tau=(-\widehat a_T-\nabla\widehat\phi_T(\widetilde y_j))_{j=1}^J,\quad\sqrt{n/L_n}(\widehat\tau-\tau(P))\Rightarrow N(0,V(P))\); the Gaussian arrow remains conditional on discharging oeq:uniform-critical-control..
\begin{itemize}
  \item \(\mathbb T^2\) --- \texttt{compact manifold}; \(\mathbb R^2/\mathbb Z^2\); \texttt{outcome space}
  \par\smallskip\noindent\emph{Definition.} The two-dimensional flat unit torus, represented by periodic lifts to \(\mathbb R^2\).
  \item \(\mathcal G\) --- \texttt{finite index set}; \(\{C0,C1,T0,T1\}\); \texttt{cell index set}
  \par\smallskip\noindent\emph{Definition.} \texttt{The control-\allowbreak{}pre,\allowbreak{} control-\allowbreak{}post,\allowbreak{} treated-\allowbreak{}pre,\allowbreak{} and treated-\allowbreak{}post cells.}
  \item \(g\) --- \texttt{index}; \(\mathcal G\); \texttt{index}
  \par\smallskip\noindent\emph{Definition.} \texttt{A generic design cell.}
  \item \(s\) --- \texttt{integer}; \(\{8,9,\ldots\}\); \texttt{smoothness index}
  \par\smallskip\noindent\emph{Definition.} \texttt{The density smoothness and kernel order.}
  \item \(\ell\) --- \texttt{index}; \(\{1,\ldots,s-1\}\); \texttt{index}
  \par\smallskip\noindent\emph{Definition.} \texttt{A kernel-\allowbreak{}moment index.}
  \item \(c\) --- \texttt{real}; \((0,1)\); \texttt{class constant}
  \par\smallskip\noindent\emph{Definition.} \texttt{The common density lower bound.}
  \item \(C\) --- \texttt{real}; \((1,\infty)\); \texttt{class constant}
  \par\smallskip\noindent\emph{Definition.} \texttt{The common density upper bound.}
  \item \(R\) --- \texttt{real}; \((0,\infty)\); \texttt{class constant}
  \par\smallskip\noindent\emph{Definition.} The radius of the periodic \(C^s\) density ball.
  \item \(\lambda\) --- \texttt{real}; \((0,1)\); \texttt{class constant}
  \par\smallskip\noindent\emph{Definition.} \texttt{The common lower ellipticity constant for the population transport Jacobians.}
  \item \(\Lambda\) --- \texttt{real}; \((1,\infty)\); \texttt{class constant}
  \par\smallskip\noindent\emph{Definition.} \texttt{The common upper ellipticity constant for the population transport Jacobians.}
  \item \(\bar\rho\) --- \texttt{real}; \([1,\infty)\); \texttt{design constant}
  \par\smallskip\noindent\emph{Definition.} \texttt{A uniform upper bound on limiting cell-\allowbreak{}size ratios.}
  \item \(\gamma\) --- \texttt{real}; \((0,1)\); \texttt{inference constant}
  \par\smallskip\noindent\emph{Definition.} \texttt{The nominal simultaneous noncoverage probability.}
  \item \(J\) --- \texttt{positive integer}; \(\mathbb N\); \texttt{grid size}
  \par\smallskip\noindent\emph{Definition.} \texttt{The fixed number of treated-\allowbreak{}post evaluation points.}
  \item \(j\) --- \texttt{index}; \(\{1,\ldots,J\}\); \texttt{index}
  \par\smallskip\noindent\emph{Definition.} \texttt{An evaluation-\allowbreak{}point index.}
  \item \(k\) --- \texttt{index}; \(\{1,\ldots,J\}\); \texttt{index}
  \par\smallskip\noindent\emph{Definition.} \texttt{A second evaluation-\allowbreak{}point index.}
  \item \(\delta_0\) --- \texttt{real}; \((0,1/2)\); \texttt{grid constant}
  \par\smallskip\noindent\emph{Definition.} \texttt{The minimum torus separation of distinct evaluation points.}
  \item \(\mathcal Y_J\) --- \texttt{ordered tuple}; \((\mathbb T^2)^J\); \texttt{evaluation grid}
  \par\smallskip\noindent\emph{Definition.} The fixed grid \(\mathcal Y_J=(y_1,\ldots,y_J)\).
  \item \(n_g\) --- \texttt{positive integer}; \(\mathbb N\); \texttt{sample size}
  \par\smallskip\noindent\emph{Definition.} The sample size in cell \(g\).
  \item \(n\) --- \texttt{positive integer}; \(\mathbb N\); \texttt{asymptotic index}
  \par\smallskip\noindent\emph{Definition.} The effective sample size \(n=\min_{g\in\mathcal G}n_g\).
  \item \(i\) --- \texttt{index}; \(\{1,\ldots,n_g\}\); \texttt{index}
  \par\smallskip\noindent\emph{Definition.} An observation index within cell \(g\).
  \item \(r\) --- \texttt{index}; \(\{1,2\}\); \texttt{index}
  \par\smallskip\noindent\emph{Definition.} \texttt{A cross-\allowbreak{}fitting fold index.}
  \item \(\rho_g\) --- \texttt{real}; \([1,\bar\rho]\); \texttt{sampling ratio}
  \par\smallskip\noindent\emph{Definition.} The limit of \(n_g/n\) for cell \(g\).
  \item \(a_h\) --- \texttt{real}; \([1/[2(s-2)],1/6)\); \texttt{smoothing regime}
  \par\smallskip\noindent\emph{Definition.} \texttt{The fixed polynomial bandwidth exponent in the maximal mixed-\allowbreak{}remainder window.}
  \item \(h_n\) --- \texttt{positive real sequence}; \((0,1)\); \texttt{bandwidth}
  \par\smallskip\noindent\emph{Definition.} The bandwidth sequence \(h_n=n^{-a_h}\).
  \item \(L_n\) --- \texttt{positive real sequence}; \((0,\infty)\); \texttt{normalization}
  \par\smallskip\noindent\emph{Definition.} The critical logarithm \(L_n=\log(1/h_n)=a_h\log n\).
  \item \(h\) --- \texttt{positive real}; \((0,1)\); \texttt{generic bandwidth}
  \par\smallskip\noindent\emph{Definition.} \texttt{A generic bandwidth used to define prelimit smoothing objects.}
  \item \(K\) --- \texttt{function}; \(C_c^\infty(\mathbb R)\); \texttt{smoothing kernel}
  \par\smallskip\noindent\emph{Definition.} \texttt{A fixed univariate signed kernel used in a separable periodized density estimator.}
  \item \(K_h^{\mathrm{per}}\) --- \texttt{function}; \(C^\infty(\mathbb T^2)\); \texttt{periodized kernel}
  \par\smallskip\noindent\emph{Definition.} The scaled periodization \(K_h^{\mathrm{per}}(x)=\sum_{k\in\mathbb Z^2}h^{-2}K((x_1+k_1)/h)K((x_2+k_2)/h)\).
  \item \(P_g\) --- \texttt{probability law}; \(\mathcal P(\mathbb T^2)\); \texttt{observed law}
  \par\smallskip\noindent\emph{Definition.} The observed law in cell \(g\).
  \item \(f_g\) --- \texttt{function}; \(C^s(\mathbb T^2)\); \texttt{cell density}
  \par\smallskip\noindent\emph{Definition.} The Lebesgue density of \(P_g\).
  \item \(X_{g,i}\) --- \texttt{random vector}; \(\mathbb T^2\); \texttt{data}
  \par\smallskip\noindent\emph{Definition.} Observation \(i\) from cell \(g\).
  \item \(P\) --- \texttt{four-\allowbreak{}law tuple}; \(\mathcal P(\mathbb T^2)^4\); \texttt{data-\allowbreak{}generating law}
  \par\smallskip\noindent\emph{Definition.} The observed-data law \(P=(P_{C0},P_{C1},P_{T0},P_{T1})\).
  \item \(c_{\mathbb T}\) --- \texttt{cost function}; \(\mathbb T^2\times\mathbb T^2\to[0,1/4]\); \texttt{torus transport cost}
  \par\smallskip\noindent\emph{Definition.} The intrinsic squared-distance cost \(c_{\mathbb T}(x,y)=d_{\mathbb T^2}(x,y)^2/2\).
  \item \(\nu_C\) --- \texttt{probability law}; \(\mathcal P(\mathbb T^2)\); \texttt{latent law}
  \par\smallskip\noindent\emph{Definition.} \texttt{The latent control rank law in a torus changes-\allowbreak{}in-\allowbreak{}changes representation.}
  \item \(\nu_T\) --- \texttt{probability law}; \(\mathcal P(\mathbb T^2)\); \texttt{latent law}
  \par\smallskip\noindent\emph{Definition.} \texttt{The latent treated rank law in a torus changes-\allowbreak{}in-\allowbreak{}changes representation.}
  \item \(h_0\) --- \texttt{diffeomorphism}; \(\mathbb T^2\to\mathbb T^2\); \texttt{production map}
  \par\smallskip\noindent\emph{Definition.} \texttt{The invertible pre-\allowbreak{}period production map in a torus changes-\allowbreak{}in-\allowbreak{}changes representation.}
  \item \(h_1\) --- \texttt{diffeomorphism}; \(\mathbb T^2\to\mathbb T^2\); \texttt{production map}
  \par\smallskip\noindent\emph{Definition.} \texttt{The invertible untreated-\allowbreak{}post production map in a torus changes-\allowbreak{}in-\allowbreak{}changes representation.}
  \item \(h_{1\star}\) --- \texttt{diffeomorphism}; \(\mathbb T^2\to\mathbb T^2\); \texttt{production map}
  \par\smallskip\noindent\emph{Definition.} \texttt{The invertible treated-\allowbreak{}post production map in a torus changes-\allowbreak{}in-\allowbreak{}changes representation.}
  \item \(\mathcal M_{\mathrm{CiC}}^{\mathrm{per}}\) --- \texttt{causal model class}; \(\mathcal P(\mathbb T^2)^4\); \texttt{new torus causal model}
  \par\smallskip\noindent\emph{Definition.} The torus changes-in-changes causal class of four observed laws admitting latent laws \(\nu_C,\nu_T\) and invertible production maps \(h_0,h_1,h_{1\star}\) with \(P_{C0}=h_0\#\nu_C\), \(P_{C1}=h_1\#\nu_C\), \(P_{T0}=h_0\#\nu_T\), \(P_{T1}=h_{1\star}\#\nu_T\), and with the graphs of \(h_1\circ h_0^{-1}\) and \(h_1\circ h_{1\star}^{-1}\) both \(c_{\mathbb T}\)-cyclically monotone.
  \item \(\operatorname{OT}\) --- \texttt{map-\allowbreak{}valued functional}; \(\mathcal P(\mathbb T^2)\times\mathcal P(\mathbb T^2)\to(\mathbb T^2\to\mathbb T^2)\); \texttt{transport operator}
  \par\smallskip\noindent\emph{Definition.} The unique almost-everywhere optimal map for the intrinsic cost \(c_{\mathbb T}\) between absolutely continuous torus laws; on the smooth class it is an identity-homotopic torus diffeomorphism.
  \item \(\mathcal Q\) --- \texttt{half-\allowbreak{}open fundamental cell}; \([-1/2,1/2)^2\); \texttt{lift branch cell}
  \par\smallskip\noindent\emph{Definition.} \texttt{The half-\allowbreak{}open cell used to select the harmonic translation of an equivariant lift.}
  \item \(F\) --- \texttt{generic map}; \(\mathbb T^2\to\mathbb T^2\); \texttt{lift argument}
  \par\smallskip\noindent\emph{Definition.} \texttt{A generic identity-\allowbreak{}homotopic torus optimal-\allowbreak{}transport diffeomorphism.}
  \item \(a_F\) --- \texttt{vector}; \(\mathcal Q\); \texttt{harmonic translation}
  \par\smallskip\noindent\emph{Definition.} The uniquely selected cell-average displacement of the equivariant lift of \(F\).
  \item \(\phi_F\) --- \texttt{periodic function}; \(C^2(\mathbb T^2)\); \texttt{lift potential}
  \par\smallskip\noindent\emph{Definition.} The mean-zero periodic potential satisfying \(\int_{\mathbb T^2}\phi_F=0\) for the selected lift of \(F\).
  \item \(\widetilde F\) --- \texttt{equivariant lift}; \(\mathbb R^2\to\mathbb R^2\); \texttt{translation-\allowbreak{}aware lift}
  \par\smallskip\noindent\emph{Definition.} The selected lift \(\widetilde F(x)=x+a_F+\nabla\phi_F(x)\), satisfying \(\widetilde F(x+k)=\widetilde F(x)+k\) for \(k\in\mathbb Z^2\).
  \item \(\operatorname{Lift}_{\mathbb T}\) --- \texttt{measurable lift selector}; \(F\mapsto(a_F,\phi_F,\widetilde F)\); \texttt{population selector and finite-\allowbreak{}sample fallback}
  \par\smallskip\noindent\emph{Definition.} For an identity-homotopic torus optimal-transport diffeomorphism, choose the unique \(\mathbb Z^2\)-equivariant representation \(\widetilde F(x)=x+a_F+\nabla\phi_F(x)\) with \(a_F\in\mathcal Q\) and \(\int_{\mathbb T^2}\phi_F=0\); outside this regular domain use the fixed Borel fallback \((a_F,\phi_F,\widetilde F)=(0,0,\operatorname{id})\).
  \item \(\#\) --- \texttt{pushforward operator}; \((\mathbb T^2\to\mathbb T^2)\times\mathcal P(\mathbb T^2)\to\mathcal P(\mathbb T^2)\); \texttt{operator}
  \par\smallskip\noindent\emph{Definition.} \texttt{The measure pushforward operation.}
  \item \(p\) --- \texttt{diffeomorphism}; \(\mathbb T^2\to\mathbb T^2\); \texttt{identified inner map}
  \par\smallskip\noindent\emph{Definition.} The inner control trend \(p=\operatorname{OT}(P_{C0},P_{C1})\).
  \item \(a_p\) --- \texttt{vector}; \(\mathcal Q\); \texttt{inner lift translation}
  \par\smallskip\noindent\emph{Definition.} The harmonic translation selected by \(\operatorname{Lift}_{\mathbb T}(p)\).
  \item \(\phi_p\) --- \texttt{periodic function}; \(C^2(\mathbb T^2)\); \texttt{inner lift potential}
  \par\smallskip\noindent\emph{Definition.} The mean-zero potential selected by \(\operatorname{Lift}_{\mathbb T}(p)\), so \(\widetilde p(x)=x+a_p+\nabla\phi_p(x)\).
  \item \(\widetilde p\) --- \texttt{equivariant lift}; \(\mathbb R^2\to\mathbb R^2\); \texttt{inner selected lift}
  \par\smallskip\noindent\emph{Definition.} The selected inner lift \(\widetilde p(x)=x+a_p+\nabla\phi_p(x)\).
  \item \(q\) --- \texttt{probability law}; \(\mathcal P(\mathbb T^2)\); \texttt{identified generated target}
  \par\smallskip\noindent\emph{Definition.} The generated treated-post counterfactual law \(q=p\#P_{T0}\).
  \item \(T\) --- \texttt{diffeomorphism}; \(\mathbb T^2\to\mathbb T^2\); \texttt{identified causal map}
  \par\smallskip\noindent\emph{Definition.} The outer counterfactual map \(T=\operatorname{OT}(P_{T1},q)\).
  \item \(a_T\) --- \texttt{vector}; \(\mathcal Q\); \texttt{outer lift translation}
  \par\smallskip\noindent\emph{Definition.} The harmonic translation selected by \(\operatorname{Lift}_{\mathbb T}(T)\).
  \item \(\phi_T\) --- \texttt{periodic function}; \(C^2(\mathbb T^2)\); \texttt{outer lift potential}
  \par\smallskip\noindent\emph{Definition.} The mean-zero potential selected by \(\operatorname{Lift}_{\mathbb T}(T)\), so \(\widetilde T(x)=x+a_T+\nabla\phi_T(x)\).
  \item \(\widetilde T\) --- \texttt{equivariant lift}; \(\mathbb R^2\to\mathbb R^2\); \texttt{outer selected lift}
  \par\smallskip\noindent\emph{Definition.} The selected outer lift \(\widetilde T(x)=x+a_T+\nabla\phi_T(x)\).
  \item \(\operatorname{Cut}_{\mathbb T}\) --- \texttt{closed subset}; \(\mathbb T^2\times\mathbb T^2\); \texttt{cost branch boundary}
  \par\smallskip\noindent\emph{Definition.} The squared-distance cut-locus relation \(\{(x,y):y-x\text{ has a coordinate congruent to }1/2\}\).
  \item \(\beta\) --- \texttt{real}; \((0,1/4)\); \texttt{lift regularity constant}
  \par\smallskip\noindent\emph{Definition.} \texttt{The fixed lift-\allowbreak{}branch and cut-\allowbreak{}locus margin.}
  \item \(\Delta_{\mathrm{lift}}(P)\) --- \texttt{nonnegative functional}; \texttt{open branch margin}
  \par\smallskip\noindent\emph{Definition.} The minimum of \(\operatorname{dist}(a_p,\partial[-1/2,1/2]^2)\), \(\operatorname{dist}(a_T,\partial[-1/2,1/2]^2)\), \(\inf_x\operatorname{dist}((x,p(x)),\operatorname{Cut}_{\mathbb T})\), and \(\inf_y\operatorname{dist}((y,T(y)),\operatorname{Cut}_{\mathbb T})\).
  \item \(\operatorname{Disp}_{\mathbb T}\) --- \texttt{displacement functional}; \((\mathbb T^2\to\mathbb T^2)\times\mathbb T^2\to\mathbb R^2\); \texttt{translation-\allowbreak{}aware torus displacement}
  \par\smallskip\noindent\emph{Definition.} If \(\operatorname{Lift}_{\mathbb T}(F)=(a_F,\phi_F,\widetilde F)\), define \(\operatorname{Disp}_{\mathbb T}(F,y)=\widetilde y-\widetilde F(\widetilde y)=-a_F-\nabla\phi_F(\widetilde y)\); equivariance makes the value independent of the representative \(\widetilde y\) of \(y\), and the fixed Borel fallback makes the functional total.
  \item \(\tau(P)\) --- \texttt{vector}; \(\mathbb R^{2J}\); \texttt{causal estimand}
  \par\smallskip\noindent\emph{Definition.} The fixed-grid translation-aware displacement \(\tau(P)=(-a_T-\nabla\phi_T(\widetilde y_j))_{j=1}^J=(\operatorname{Disp}_{\mathbb T}(T,y_j))_{j=1}^J\).
  \item \(R_0\) --- \texttt{diffeomorphism}; \(\mathbb T^2\to\mathbb T^2\); \texttt{coordinate pullback}
  \par\smallskip\noindent\emph{Definition.} The treated-post to control-pre coordinate map \(R_0=p^{-1}\circ T\).
  \item \(\alpha\) --- \texttt{function}; \(C^{s-1}(\mathbb T^2)\); \texttt{nested noise weight}
  \par\smallskip\noindent\emph{Definition.} The pulled-back density ratio \(\alpha=(f_{T0}/f_{C0})\circ R_0\).
  \item \(S\) --- \texttt{matrix-\allowbreak{}valued function}; \(\mathbb T^2\to\mathbb R^{2\times2}\); \texttt{inner coefficient}
  \par\smallskip\noindent\emph{Definition.} The inner Jacobian \(S=Dp\).
  \item \(H\) --- \texttt{matrix-\allowbreak{}valued function}; \(\mathbb T^2\to\mathbb R^{2\times2}\); \texttt{outer coefficient}
  \par\smallskip\noindent\emph{Definition.} The outer Jacobian \(H=DT\).
  \item \(D\) --- \texttt{matrix-\allowbreak{}valued function}; \(\mathbb T^2\to\mathbb R^{2\times2}\); \texttt{coordinate Jacobian}
  \par\smallskip\noindent\emph{Definition.} The Jacobian \(D=DR_0=S(R_0)^{-1}H\).
  \item \(A\) --- \texttt{matrix-\allowbreak{}valued function}; \(\mathbb T^2\to\mathbb R^{2\times2}\); \texttt{elliptic coefficient}
  \par\smallskip\noindent\emph{Definition.} The inner elliptic coefficient \(A=f_{C0}S^{-1}\).
  \item \(B\) --- \texttt{matrix-\allowbreak{}valued function}; \(\mathbb T^2\to\mathbb R^{2\times2}\); \texttt{elliptic coefficient}
  \par\smallskip\noindent\emph{Definition.} The outer elliptic coefficient \(B=f_{T1}H^{-1}\).
  \item \(\mathsf C\) --- \texttt{matrix-\allowbreak{}valued function}; \(\mathbb T^2\to\mathbb R^{2\times2}\); \texttt{nested elliptic coefficient}
  \par\smallskip\noindent\emph{Definition.} The pulled-back inner coefficient \(\mathsf C=\det(D)D^{-1}A(R_0)D^{-\mathsf T}\).
  \item \(\Pi_0\) --- \texttt{linear operator}; \(u\mapsto u-\int_{\mathbb T^2}u\); \texttt{normalization}
  \par\smallskip\noindent\emph{Definition.} \texttt{Projection onto mean-\allowbreak{}zero periodic functions.}
  \item \(\mathcal L_A\) --- \texttt{elliptic operator}; \(u\mapsto-\operatorname{div}(A\nabla u)\); \texttt{linearization operator}
  \par\smallskip\noindent\emph{Definition.} \texttt{The normalized inner linearized Monge-\allowbreak{}-\allowbreak{}Ampere operator.}
  \item \(\mathcal L_B\) --- \texttt{elliptic operator}; \(u\mapsto-\operatorname{div}(B\nabla u)\); \texttt{linearization operator}
  \par\smallskip\noindent\emph{Definition.} \texttt{The normalized outer linearized Monge-\allowbreak{}-\allowbreak{}Ampere operator.}
  \item \(\mathcal L_{\mathsf C}\) --- \texttt{elliptic operator}; \(u\mapsto-\operatorname{div}(\mathsf C\nabla u)\); \texttt{linearization operator}
  \par\smallskip\noindent\emph{Definition.} \texttt{The normalized pulled-\allowbreak{}back inner elliptic operator.}
  \item \(\delta\) --- \texttt{density tangent}; \((C_0^{s-2}(\mathbb T^2))^4\); \texttt{perturbation direction}
  \par\smallskip\noindent\emph{Definition.} A four-cell mean-zero tangent \(\delta=(\delta_{C0},\delta_{C1},\delta_{T0},\delta_{T1})\).
  \item \(e_{C0}\) --- \texttt{function}; \(C_0^{s-2}(\mathbb T^2)\); \texttt{transformed noise}
  \par\smallskip\noindent\emph{Definition.} The pullback \(e_{C0}=\det(D)\,\delta_{C0}\circ R_0\).
  \item \(e_{C1}\) --- \texttt{function}; \(C_0^{s-2}(\mathbb T^2)\); \texttt{transformed noise}
  \par\smallskip\noindent\emph{Definition.} The pullback \(e_{C1}=\det(H)\,\delta_{C1}\circ T\).
  \item \(e_{T0}\) --- \texttt{function}; \(C_0^{s-2}(\mathbb T^2)\); \texttt{transformed noise}
  \par\smallskip\noindent\emph{Definition.} The pullback \(e_{T0}=\det(D)\,\delta_{T0}\circ R_0\).
  \item \(e_{T1}\) --- \texttt{function}; \(C_0^{s-2}(\mathbb T^2)\); \texttt{transformed noise}
  \par\smallskip\noindent\emph{Definition.} The treated-post tangent \(e_{T1}=\delta_{T1}\).
  \item \(\mathcal H_P\) --- \texttt{linear operator}; \((C_0^{s-2}(\mathbb T^2))^4\to(\mathbb R^2)^J\); \texttt{total coupled harmonic correction}
  \par\smallskip\noindent\emph{Definition.} \texttt{The complete correction to the zero-\allowbreak{}harmonic pulled-\allowbreak{}coordinate Green response,\allowbreak{} including the inner and outer harmonic increments and the scalar elliptic responses generated by their coefficient-\allowbreak{}divergence terms.}
  \item \(\mathcal D_P\) --- \texttt{linear operator}; \((C_0^{s-2}(\mathbb T^2))^4\to(\mathbb R^2)^J\); \texttt{nested pathwise derivative}
  \par\smallskip\noindent\emph{Definition.} \texttt{The coupled scalar-\allowbreak{}-\allowbreak{}harmonic derivative of the selected Euclidean lifts,\allowbreak{} obtained by differentiating both mass equations and both finite-\allowbreak{}dimensional reduced-\allowbreak{}cost optimality equations and propagating the inner field through the generated target law.}
  \item \(\widetilde f_g\) --- \texttt{random function}; \(C^\infty(\mathbb T^2)\); \texttt{density pilot}
  \par\smallskip\noindent\emph{Definition.} The signed KDE \(\widetilde f_g(x)=n_g^{-1}\sum_{i=1}^{n_g}K_{h_n}^{\mathrm{per}}(x-X_{g,i})\).
  \item \(\widehat f_g\) --- \texttt{random density}; \(C(\mathbb T^2)\); \texttt{density estimator}
  \par\smallskip\noindent\emph{Definition.} The positive normalized density \(\widehat f_g(x)=\max\{c/2,\widetilde f_g(x)\}/\int_{\mathbb T^2}\max\{c/2,\widetilde f_g(u)\}\,du\).
  \item \(\widehat P_g\) --- \texttt{random probability law}; \(\mathcal P(\mathbb T^2)\); \texttt{cell-\allowbreak{}law estimator}
  \par\smallskip\noindent\emph{Definition.} The law with density \(\widehat f_g\).
  \item \(\widehat p\) --- \texttt{random map}; \(\mathbb T^2\to\mathbb T^2\); \texttt{inner estimator}
  \par\smallskip\noindent\emph{Definition.} The inner plug-in map \(\widehat p=\operatorname{OT}(\widehat P_{C0},\widehat P_{C1})\), where \(\widehat P_g\) has density \(\widehat f_g\).
  \item \(\widehat a_p\) --- \texttt{random vector}; \(\mathcal Q\); \texttt{estimated inner translation}
  \par\smallskip\noindent\emph{Definition.} The half-open-cell translation returned by \(\operatorname{Lift}_{\mathbb T}(\widehat p)\), with its measurable fallback on the exceptional event.
  \item \(\widehat\phi_p\) --- \texttt{random periodic function}; \(C^2(\mathbb T^2)\); \texttt{estimated inner potential}
  \par\smallskip\noindent\emph{Definition.} The mean-zero potential returned by \(\operatorname{Lift}_{\mathbb T}(\widehat p)\), with its measurable fallback on the exceptional event.
  \item \(\widetilde{\widehat p}\) --- \texttt{random equivariant lift}; \(\mathbb R^2\to\mathbb R^2\); \texttt{estimated inner selected lift}
  \par\smallskip\noindent\emph{Definition.} The lift returned with \(\widehat a_p\) and \(\widehat\phi_p\) by \(\operatorname{Lift}_{\mathbb T}(\widehat p)\), including its measurable fallback.
  \item \(\widehat q\) --- \texttt{random probability law}; \(\mathcal P(\mathbb T^2)\); \texttt{generated target estimator}
  \par\smallskip\noindent\emph{Definition.} The generated plug-in law \(\widehat q=\widehat p\#\widehat P_{T0}\), evaluated through \(\widehat f_q(z)=\widehat f_{C1}(z)(\widehat f_{T0}/\widehat f_{C0})(\widehat p^{-1}(z))\).
  \item \(\widehat T\) --- \texttt{random map}; \(\mathbb T^2\to\mathbb T^2\); \texttt{map estimator}
  \par\smallskip\noindent\emph{Definition.} The outer plug-in map \(\widehat T=\operatorname{OT}(\widehat P_{T1},\widehat q)\).
  \item \(\widehat a_T\) --- \texttt{random vector}; \(\mathcal Q\); \texttt{estimated outer translation}
  \par\smallskip\noindent\emph{Definition.} The half-open-cell translation returned by \(\operatorname{Lift}_{\mathbb T}(\widehat T)\), with its measurable fallback on the exceptional event.
  \item \(\widehat\phi_T\) --- \texttt{random periodic function}; \(C^2(\mathbb T^2)\); \texttt{estimated outer potential}
  \par\smallskip\noindent\emph{Definition.} The mean-zero potential returned by \(\operatorname{Lift}_{\mathbb T}(\widehat T)\), with its measurable fallback on the exceptional event.
  \item \(\widetilde{\widehat T}\) --- \texttt{random equivariant lift}; \(\mathbb R^2\to\mathbb R^2\); \texttt{estimated outer selected lift}
  \par\smallskip\noindent\emph{Definition.} The lift returned with \(\widehat a_T\) and \(\widehat\phi_T\) by \(\operatorname{Lift}_{\mathbb T}(\widehat T)\), including its measurable fallback.
  \item \(\widehat\tau\) --- \texttt{random vector}; \(\mathbb R^{2J}\); \texttt{causal estimator}
  \par\smallskip\noindent\emph{Definition.} The total translation-aware displacement estimator \(\widehat\tau=(-\widehat a_T-\nabla\widehat\phi_T(\widetilde y_j))_{j=1}^J=(\operatorname{Disp}_{\mathbb T}(\widehat T,y_j))_{j=1}^J\), using the measurable fallback returned by the lift selector whenever its regular branch is unavailable.
  \item \(\Psi_{g,h}\) --- \texttt{vector-\allowbreak{}valued function}; \(L^2(P_g;\mathbb R^{2J})\); \texttt{finite-\allowbreak{}bandwidth score}
  \par\smallskip\noindent\emph{Definition.} The centered full nested influence vector obtained by applying the exact derivative \(\mathcal D_P\), including \(\mathcal H_P\), to the cell-\(g\) tangent \(K_h^{\mathrm{per}}(\cdot-x)-K_h^{\mathrm{per}}*f_g\), setting the other three tangents to zero, and stacking the evaluations on \(\mathcal Y_J\).
  \item \(\Sigma_h(P)\) --- \texttt{matrix}; \(\mathbb S_+^{2J}\); \texttt{prelimit covariance}
  \par\smallskip\noindent\emph{Definition.} The full normalized score covariance \(\Sigma_h(P)=\log(1/h)^{-1}\sum_{g\in\mathcal G}\rho_g^{-1}\operatorname{Var}_{P_g}(\Psi_{g,h}(X))\).
  \item \(V(P)\) --- \texttt{matrix}; \(\mathbb S_{++}^{2J}\); \texttt{limiting covariance}
  \par\smallskip\noindent\emph{Definition.} The block-diagonal critical covariance \(V(P)=\operatorname{diag}(V_1,\ldots,V_J)\), with \(V_j=w_j\sqrt{\det H(y_j)}H(y_j)/(4\pi f_{T1}(y_j))\) and \(w_j=\alpha(y_j)(\rho_{C0}^{-1}+\rho_{C1}^{-1})+\rho_{T0}^{-1}+\rho_{T1}^{-1}\).
  \item \(\widehat\Psi_{g,h}^{(-r)}\) --- \texttt{vector-\allowbreak{}valued random function}; \(L^2(P_g;\mathbb R^{2J})\); \texttt{cross-\allowbreak{}fitted adjoint score}
  \par\smallskip\noindent\emph{Definition.} The score \(\Psi_{g,h}\) computed by adjoint solves with density, coefficient, map, and pole pilots from the fold opposite evaluation fold \(r\in\{1,2\}\).
  \item \(\widehat V_{\mathrm{adj}}\) --- \texttt{random matrix}; \(\mathbb S_+^{2J}\); \texttt{covariance estimator}
  \par\smallskip\noindent\emph{Definition.} The two-fold cross-fitted empirical covariance \(\widehat V_{\mathrm{adj}}=L_n^{-1}\sum_{g\in\mathcal G}(n/n_g)n_g^{-1}\sum_{i=1}^{n_g}(\widehat\Psi_{g,h_n}^{(-r(i))}(X_{g,i})-\overline\Psi_{g,r(i)})(\widehat\Psi_{g,h_n}^{(-r(i))}(X_{g,i})-\overline\Psi_{g,r(i)})^{\mathsf T}\), where \(\overline\Psi_{g,r}\) is the score mean within evaluation fold \(r\).
  \item \(\theta_n\) --- \texttt{positive deterministic sequence}; \((0,\infty)\); \texttt{finite-\allowbreak{}sample regularization}
  \par\smallskip\noindent\emph{Definition.} The ridge sequence \(\theta_n=n^{-1}\).
  \item \(\widehat V_{\mathrm{adj}}^{R}\) --- \texttt{random matrix}; \(\mathbb S_{++}^{2J}\); \texttt{ridged covariance estimator}
  \par\smallskip\noindent\emph{Definition.} The totalized covariance matrix \(\widehat V_{\mathrm{adj}}^{R}=\widehat V_{\mathrm{adj}}+\theta_n I_{2J}\).
  \item \(\xi_{g,i}\) --- \texttt{random variable}; \(N(0,1)\); \texttt{multiplier}
  \par\smallskip\noindent\emph{Definition.} \texttt{An independent Gaussian multiplier,\allowbreak{} conditionally independent of the data.}
  \item \(Z^\star\) --- \texttt{random vector}; \(\mathbb R^{2J}\); \texttt{calibration statistic}
  \par\smallskip\noindent\emph{Definition.} The conditional multiplier vector \(Z^\star=\sqrt{n/L_n}\sum_{g\in\mathcal G}n_g^{-1}\sum_{i=1}^{n_g}\xi_{g,i}(\widehat\Psi_{g,h_n}^{(-r(i))}(X_{g,i})-\overline\Psi_{g,r(i)})\).
  \item \(\widehat c_{1-\gamma}\) --- \texttt{random scalar}; \([0,\infty)\); \texttt{simultaneous critical value}
  \par\smallskip\noindent\emph{Definition.} The conditional \(1-\gamma\) quantile, allowed to equal zero, of \(\max_{1\le j\le J}(Z_j^\star)^{\mathsf T}(\widehat V_{\mathrm{adj},jj}^{R})^{-1}Z_j^\star\).
  \item \(\mathcal E_j\) --- \texttt{random set}; \(2^{\mathbb R^2}\); \texttt{confidence set}
  \par\smallskip\noindent\emph{Definition.} The total finite-sample ellipse \(\mathcal E_j=\{u\in\mathbb R^2:(n/L_n)(\widehat\tau_j-u)^{\mathsf T}(\widehat V_{\mathrm{adj},jj}^{R})^{-1}(\widehat\tau_j-u)\le\widehat c_{1-\gamma}\}\).
  \item \(P_0\) --- \texttt{four-\allowbreak{}law tuple}; \(\mathcal P(\mathbb T^2)^4\); \texttt{local minimax center}
  \par\smallskip\noindent\emph{Definition.} \texttt{A four-\allowbreak{}law tuple used as the center of the local experiment.}
  \item \(P_{0g}\) --- \texttt{probability law}; \(\mathcal P(\mathbb T^2)\); \texttt{local-\allowbreak{}center cell law}
  \par\smallskip\noindent\emph{Definition.} The cell-\(g\) component of the local center \(P_0\).
  \item \(f_{0g}\) --- \texttt{function}; \(C^s(\mathbb T^2)\); \texttt{local-\allowbreak{}center density}
  \par\smallskip\noindent\emph{Definition.} The density of cell \(g\) in the local center \(P_0\).
  \item \(\kappa\) --- \texttt{real}; \((0,\infty)\); \texttt{local minimax constant}
  \par\smallskip\noindent\emph{Definition.} \texttt{The product-\allowbreak{}Hellinger neighborhood radius.}
  \item \(\mathcal F_s\) --- \texttt{model class}; \(2^{\mathcal P(\mathbb T^2)^4}\); \texttt{law class}
  \par\smallskip\noindent\emph{Definition.} The bounded smooth periodic submodel of \(\mathcal M_{\mathrm{CiC}}^{\mathrm{per}}\).
  \item \(\operatorname{Hel}\) --- \texttt{statistical distance}; \texttt{local-\allowbreak{}experiment metric}
  \par\smallskip\noindent\emph{Definition.} \texttt{The Hellinger distance between probability laws on a common measurable space.}
  \item \(\mathcal N_n(P_0,\kappa)\) --- \texttt{set of four-\allowbreak{}law tuples}; \(2^{\mathcal P(\mathbb T^2)^4}\); \texttt{local experiment}
  \par\smallskip\noindent\emph{Definition.} The local neighborhood \(\mathcal N_n(P_0,\kappa)=\{P\in\mathcal F_s:\sum_{g\in\mathcal G}\operatorname{Hel}^2(P_g,P_{0g})\le\kappa/n\}\).
  \item \(g_\star\) --- \texttt{cell index sequence}; \(\mathcal G\); \texttt{local-\allowbreak{}converse perturbation cell}
  \par\smallskip\noindent\emph{Definition.} A deterministically tie-broken least-sampled cell satisfying \(n_{g_\star}=n\).
  \item \(\eta\) --- \texttt{real}; \((0,1/[2(s+1)])\); \texttt{local-\allowbreak{}converse scale}
  \par\smallskip\noindent\emph{Definition.} \texttt{The polynomial inner-\allowbreak{}cutoff exponent for the local alternatives.}
  \item \(\varepsilon\) --- \texttt{positive real}; \((0,\delta_0/8)\); \texttt{local-\allowbreak{}converse generic scale}
  \par\smallskip\noindent\emph{Definition.} \texttt{A generic inner cutoff radius for the truncated Green score.}
  \item \(\varepsilon_n\) --- \texttt{positive real sequence}; \((0,1)\); \texttt{local-\allowbreak{}converse cutoff}
  \par\smallskip\noindent\emph{Definition.} The Green cutoff \(\varepsilon_n=n^{-\eta}\), chosen independently of the estimator bandwidth.
  \item \(\upsilon\) --- \texttt{vector}; \(\mathbb S^1\); \texttt{local-\allowbreak{}converse direction}
  \par\smallskip\noindent\emph{Definition.} A fixed unit direction used to scalarize the displacement at \(y_1\).
  \item \(\varphi_{g_\star}\) --- \texttt{locally square-\allowbreak{}integrable function}; \(L^2_{\mathrm{loc}}(P_{0g_\star})\); \texttt{Green-\allowbreak{}aligned local score}
  \par\smallskip\noindent\emph{Definition.} The scalar Riesz representer obtained by applying the leading outer-Green term of \(\mathcal D_{P_0}\) to a density tangent in cell \(g_\star\), evaluating at \(y_1\), and projecting the resulting displacement on \(\upsilon\).
  \item \(m\) --- \texttt{nonnegative integer}; \(\mathbb N_0\); \texttt{local-\allowbreak{}converse regularity index}
  \par\smallskip\noindent\emph{Definition.} \texttt{A derivative order for the annular cutoff bounds.}
  \item \(\chi_\varepsilon\) --- \texttt{smooth function family}; \(C^\infty(\mathbb T^2;[0,1])\); \texttt{multiscale truncation}
  \par\smallskip\noindent\emph{Definition.} A fixed annular cutoff around the pole of \(\varphi_{g_\star}\), equal to zero at torus distance at most \(\varepsilon\), equal to one from distance \(2\varepsilon\) through \(\delta_0/4\), and equal to zero from distance \(\delta_0/2\) onward, with transition derivatives of order \(m\) bounded by constants times \(\varepsilon^{-m}\) at the inner boundary.
  \item \(\ell_{\varepsilon,g_\star}\) --- \texttt{centered function}; \(L^2_0(P_{0g_\star})\); \texttt{multiscale tangent}
  \par\smallskip\noindent\emph{Definition.} The centered truncated Green score \(\ell_{\varepsilon,g_\star}=\chi_\varepsilon\varphi_{g_\star}-\int\chi_\varepsilon\varphi_{g_\star}\,dP_{0g_\star}\).
  \item \(\zeta\) --- \texttt{positive real}; \((0,1)\); \texttt{local-\allowbreak{}converse amplitude}
  \par\smallskip\noindent\emph{Definition.} \texttt{A fixed sufficiently small amplitude for the two local alternatives.}
  \item \(f_{n,\pm,g}\) --- \texttt{density family}; \(C^s(\mathbb T^2)\); \texttt{local-\allowbreak{}converse density}
  \par\smallskip\noindent\emph{Definition.} \texttt{The cell densities of the two Green-\allowbreak{}aligned alternatives constructed in def:\allowbreak{}green-\allowbreak{}local-\allowbreak{}alternatives.}
  \item \(P_{n,\pm,g}\) --- \texttt{probability-\allowbreak{}law family}; \(\mathcal P(\mathbb T^2)\); \texttt{local-\allowbreak{}converse cell law}
  \par\smallskip\noindent\emph{Definition.} The cell laws with densities \(f_{n,\pm,g}\).
  \item \(P_{n,+}\) --- \texttt{four-\allowbreak{}law tuple}; \(\mathcal P(\mathbb T^2)^4\); \texttt{local-\allowbreak{}converse alternative}
  \par\smallskip\noindent\emph{Definition.} \texttt{The positive Green-\allowbreak{}aligned alternative constructed in def:\allowbreak{}green-\allowbreak{}local-\allowbreak{}alternatives.}
  \item \(P_{n,-}\) --- \texttt{four-\allowbreak{}law tuple}; \(\mathcal P(\mathbb T^2)^4\); \texttt{local-\allowbreak{}converse alternative}
  \par\smallskip\noindent\emph{Definition.} \texttt{The negative Green-\allowbreak{}aligned alternative constructed in def:\allowbreak{}green-\allowbreak{}local-\allowbreak{}alternatives.}
  \item \(\mathfrak H_{\mathrm{crit}}\) --- \texttt{analytic proof handle}; \texttt{open-\allowbreak{}question handle}
  \par\smallskip\noindent\emph{Definition.} \texttt{The pulled-\allowbreak{}coordinate,\allowbreak{} pole-\allowbreak{}parametrix,\allowbreak{} and tame-\allowbreak{}remainder program for the nested critical response.}
  \item \(U\) --- \texttt{random vector}; \(\mathbb T^2\); \texttt{witness latent variable}
  \par\smallskip\noindent\emph{Definition.} \texttt{A uniform torus rank used in the explicit separable witness.}
  \item \(h_0^{\mathrm{sep}}\) --- \texttt{diffeomorphism}; \(\mathbb T^2\to\mathbb T^2\); \texttt{witness production map}
  \par\smallskip\noindent\emph{Definition.} The witness map \(h_0^{\mathrm{sep}}(u)=u\).
  \item \(h_1^{\mathrm{sep}}\) --- \texttt{diffeomorphism}; \(\mathbb T^2\to\mathbb T^2\); \texttt{witness untreated-\allowbreak{}post map}
  \par\smallskip\noindent\emph{Definition.} The witness map \(h_1^{\mathrm{sep}}(u)=(u_1+0.1\sin(2\pi u_1)/\pi,u_2)\).
  \item \(h_{1\star}^{\mathrm{sep}}\) --- \texttt{diffeomorphism}; \(\mathbb T^2\to\mathbb T^2\); \texttt{witness treated-\allowbreak{}post map}
  \par\smallskip\noindent\emph{Definition.} The witness map \(h_{1\star}^{\mathrm{sep}}(u)=(u_1+0.15\sin(2\pi u_1)/\pi,u_2-0.05\sin(2\pi u_2)/\pi)\).
  \item \(P^{\mathrm{sep}}\) --- \texttt{four-\allowbreak{}law tuple}; \(\mathcal P(\mathbb T^2)^4\); \texttt{explicit witness law}
  \par\smallskip\noindent\emph{Definition.} The witness laws induced by \(Y_{C0}=h_0^{\mathrm{sep}}(U)\), \(Y_{C1}=h_1^{\mathrm{sep}}(U)\), \(Y_{T0}=h_0^{\mathrm{sep}}(U)\), and \(Y_{T1}=h_{1\star}^{\mathrm{sep}}(U)\).
\end{itemize}

\section{Assumptions}

\begin{assumption}[ass:cic-identification]\label{ass:cic-identification}
There exist \(\nu_C,\nu_T,h_0,h_1,h_{1\star}\) such that \(h_0,h_1,h_{1\star}\) are invertible torus diffeomorphisms, \(P_{C0}=h_0\#\nu_C\), \(P_{C1}=h_1\#\nu_C\), \(P_{T0}=h_0\#\nu_T\), \(P_{T1}=h_{1\star}\#\nu_T\), and the graphs of \(h_1\circ h_0^{-1}\) and \(h_1\circ h_{1\star}^{-1}\) are \(c_{\mathbb T}\)-cyclically monotone.
\par\smallskip\noindent\emph{Novel.} This intrinsic flat-torus production model is newly maintained rather than inherited from the Euclidean changes-in-changes theorem. It is realistic for genuinely periodic bivariate outcomes, including seasonal phase pairs and directional coordinates, and it is satisfied by smooth equivariant gradient production maps with positive Jacobian margins.
\end{assumption}

\begin{assumption}[ass:density-positivity]\label{ass:density-positivity}
\(c\leq\min_{g\in\mathcal G}\inf_{x\in\mathbb T^2}f_g(x)\leq\max_{g\in\mathcal G}\sup_{x\in\mathbb T^2}f_g(x)\leq C\).
\par\smallskip\noindent\emph{Standard} (Uniformly positive and bounded torus densities, \cite{ManoleBalakrishnanNilesWeedWasserman2023CLT}).
\end{assumption}

\begin{assumption}[ass:density-smoothness]\label{ass:density-smoothness}
\(\max_{g\in\mathcal G}\lVert f_g\rVert_{C^s(\mathbb T^2)}\leq R\).
\par\smallskip\noindent\emph{Standard} (Periodic Holder smoothness ball, \cite{ManoleBalakrishnanNilesWeedWasserman2023CLT}).
\end{assumption}

\begin{assumption}[ass:inner-ellipticity]\label{ass:inner-ellipticity}
\(\lambda I_2\preceq S(x)\preceq\Lambda I_2\) for every \(x\in\mathbb T^2\).
\par\smallskip\noindent\emph{Standard} (Uniform ellipticity of the Brenier-map Jacobian, \cite{HutterRigollet2021Minimax}).
\end{assumption}

\begin{assumption}[ass:outer-ellipticity]\label{ass:outer-ellipticity}
\(\lambda I_2\preceq H(y)\preceq\Lambda I_2\) for every \(y\in\mathbb T^2\).
\par\smallskip\noindent\emph{Standard} (Uniform ellipticity of the Brenier-map Jacobian, \cite{HutterRigollet2021Minimax}).
\end{assumption}

\begin{assumption}[ass:lift-regularity]\label{ass:lift-regularity}
\(\Delta_{\mathrm{lift}}(P)>\beta\).
\par\smallskip\noindent\emph{Novel.} This open separation condition excludes only maps whose selected harmonic translation lies on a lattice branch boundary or whose transport graph touches the squared-distance cut locus. It is stable under small \(C^1\) perturbations and is satisfied by smooth periodic displacement fields whose coordinate displacements stay strictly inside one injectivity cell, including the explicit separable witness.
\end{assumption}

\begin{assumption}[ass:four-sample-design]\label{ass:four-sample-design}
\(\{X_{g,i}:g\in\mathcal G,1\leq i\leq n_g\}\sim\bigotimes_{g\in\mathcal G}P_g^{\otimes n_g}\).
\par\smallskip\noindent\emph{Standard} (Independent repeated-cross-section sampling, \cite{TorousGunsiliusRigollet2024OTCiC}).
\end{assumption}

\begin{assumption}[ass:cell-ratios]\label{ass:cell-ratios}
\((n_g/n)_{g\in\mathcal G}\to(\rho_g)_{g\in\mathcal G}\in[1,\bar\rho]^4\).
\par\smallskip\noindent\emph{Standard} (Stable multi-sample proportions, \cite{PonnopratOkanoImaizumi2024UniformBand}).
\end{assumption}

\begin{assumption}[ass:fixed-separated-grid]\label{ass:fixed-separated-grid}
\(\min_{j\neq k}d_{\mathbb T^2}(y_j,y_k)\geq\delta_0\).
\par\smallskip\noindent\emph{Standard} (Fixed separated evaluation grid, \cite{ManoleBalakrishnanNilesWeedWasserman2023CLT}).
\end{assumption}

\begin{assumption}[ass:kernel-symmetry]\label{ass:kernel-symmetry}
\(K(u)=K(-u)\) for every \(u\in\mathbb R\).
\par\smallskip\noindent\emph{Standard} (Symmetric smooth kernel condition, \cite{ManoleBalakrishnanNilesWeedWasserman2023CLT}).
\end{assumption}

\begin{assumption}[ass:kernel-mass]\label{ass:kernel-mass}
\(\int_{\mathbb R}K(u)\,du=1\).
\par\smallskip\noindent\emph{Standard} (Unit-mass kernel condition, \cite{ManoleBalakrishnanNilesWeedWasserman2023CLT}).
\end{assumption}

\begin{assumption}[ass:kernel-order]\label{ass:kernel-order}
\(\int_{\mathbb R}u^\ell K(u)\,du=0\) for every \(\ell\in\{1,\ldots,s-1\}\).
\par\smallskip\noindent\emph{Standard} (Order-s signed kernel moment condition, \cite{ManoleBalakrishnanNilesWeedWasserman2023CLT}).
\end{assumption}

\begin{assumption}[ass:polynomial-bandwidth]\label{ass:polynomial-bandwidth}
\(h_n=n^{-a_h}\) with \(1/[2(s-2)]\leq a_h<1/6\).
\par\smallskip\noindent\emph{Novel.} This synchronized window balances the order-minus-two bias and the mixed two-endpoint nonlinear remainder. For each bandwidth exponent below the open upper endpoint, a fixed positive Hölder index can be selected so that pilot consistency, moving-pole stability, and the normalized mixed remainder hold; at the upper endpoint the stochastic--stochastic product diverges for every positive Hölder index.
\end{assumption}

\begin{assumption}[ass:local-center]\label{ass:local-center}
\(P_0\in\operatorname{int}(\mathcal F_s)\) relative to the \(C^s\) topology.
\par\smallskip\noindent\emph{Standard} (Interior-point condition for local minimax analysis, \cite{HutterRigollet2021Minimax}).
\end{assumption}

\section{Definitions}

\begin{definition}[def:torus-cic-causal-class]\label{def:torus-cic-causal-class}
\par\smallskip\noindent\emph{Defined object.} \(\mathcal M_{\mathrm{CiC}}^{\mathrm{per}}\)
\par\smallskip\noindent\emph{Construction.} \(\mathcal M_{\mathrm{CiC}}^{\mathrm{per}}=\{P=(P_g)_{g\in\mathcal G}:\text{there exist }\nu_C,\nu_T,h_0,h_1,h_{1\star}\text{ satisfying ass:cic-identification}\}\).
\par\smallskip\noindent\emph{Member properties.} \texttt{ass:\allowbreak{}cic-\allowbreak{}identification}.
\end{definition}

\begin{definition}[def:model-class]\label{def:model-class}
\par\smallskip\noindent\emph{Defined object.} \(\mathcal F_s\)
\par\smallskip\noindent\emph{Construction.} \(\mathcal F_s=\{P=(P_g)_{g\in\mathcal G}:P\in\mathcal M_{\mathrm{CiC}}^{\mathrm{per}}\text{ and ass:density-positivity, ass:density-smoothness, ass:inner-ellipticity, ass:outer-ellipticity, and ass:lift-regularity hold}\}\).
\par\smallskip\noindent\emph{Member properties.} \texttt{ass:\allowbreak{}cic-\allowbreak{}identification}; \texttt{ass:\allowbreak{}density-\allowbreak{}positivity}; \texttt{ass:\allowbreak{}density-\allowbreak{}smoothness}; \texttt{ass:\allowbreak{}inner-\allowbreak{}ellipticity}; \texttt{ass:\allowbreak{}outer-\allowbreak{}ellipticity}; \texttt{ass:\allowbreak{}lift-\allowbreak{}regularity}.
\end{definition}

\begin{definition}[def:translation-aware-lift]\label{def:translation-aware-lift}
\par\smallskip\noindent\emph{Defined object.} \texttt{Measurable translation-\allowbreak{}aware equivariant lift}
\par\smallskip\noindent\emph{Construction.} For each regular identity-homotopic torus optimal-transport diffeomorphism \(F\), choose the unique triple \(\operatorname{Lift}_{\mathbb T}(F)=(a_F,\phi_F,\widetilde F)\) such that \(a_F\in\mathcal Q\), \(\int_{\mathbb T^2}\phi_F=0\), and \(\widetilde F(x)=x+a_F+\nabla\phi_F(x)\) is \(\mathbb Z^2\)-equivariant. Extend this selector measurably to all map estimates by the fixed fallback \((0,0,\operatorname{id})\) off the regular domain, and define \(\operatorname{Disp}_{\mathbb T}(F,y)=-a_F-\nabla\phi_F(\widetilde y)\).
\end{definition}

\begin{definition}[def:identified-functional]\label{def:identified-functional}
\par\smallskip\noindent\emph{Defined object.} \texttt{Nested Brenier changes-\allowbreak{}in-\allowbreak{}changes functional}
\par\smallskip\noindent\emph{Construction.} \(p=\operatorname{OT}(P_{C0},P_{C1}),\ q=p\#P_{T0},\ T=\operatorname{OT}(P_{T1},q)\). Select \(\operatorname{Lift}_{\mathbb T}(p)=(a_p,\phi_p,\widetilde p)\) and \(\operatorname{Lift}_{\mathbb T}(T)=(a_T,\phi_T,\widetilde T)\), and set \(\tau(P)=(-a_T-\nabla\phi_T(\widetilde y_j))_{j=1}^J=(\operatorname{Disp}_{\mathbb T}(T,y_j))_{j=1}^J\).
\par\smallskip\noindent\emph{Inputs.} \texttt{def:\allowbreak{}translation-\allowbreak{}aware-\allowbreak{}lift}.
\end{definition}

\begin{definition}[def:kde]\label{def:kde}
\par\smallskip\noindent\emph{Defined object.} \texttt{Order-\allowbreak{}s positive-\allowbreak{}normalized periodized KDE}
\par\smallskip\noindent\emph{Construction.} For every \(g\in\mathcal G\), set \(\widetilde f_g(x)=n_g^{-1}\sum_{i=1}^{n_g}K_{h_n}^{\mathrm{per}}(x-X_{g,i})\) and \(\widehat f_g(x)=\max\{c/2,\widetilde f_g(x)\}/\int_{\mathbb T^2}\max\{c/2,\widetilde f_g(u)\}\,du\).
\par\smallskip\noindent\emph{Inputs.} \texttt{ass:\allowbreak{}kernel-\allowbreak{}symmetry}; \texttt{ass:\allowbreak{}kernel-\allowbreak{}mass}; \texttt{ass:\allowbreak{}kernel-\allowbreak{}order}.
\end{definition}

\begin{definition}[def:plugin-estimator]\label{def:plugin-estimator}
\par\smallskip\noindent\emph{Defined object.} \texttt{Nested KDE and Monge-\allowbreak{}-\allowbreak{}Ampere plug-\allowbreak{}in}
\par\smallskip\noindent\emph{Construction.} Compute \(\widehat p=\operatorname{OT}(\widehat P_{C0},\widehat P_{C1})\), select \(\operatorname{Lift}_{\mathbb T}(\widehat p)=(\widehat a_p,\widehat\phi_p,\widetilde{\widehat p})\), form \(\widehat q=\widehat p\#\widehat P_{T0}\) through \(\widehat f_q(z)=\widehat f_{C1}(z)(\widehat f_{T0}/\widehat f_{C0})(\widehat p^{-1}(z))\), compute \(\widehat T=\operatorname{OT}(\widehat P_{T1},\widehat q)\), select \(\operatorname{Lift}_{\mathbb T}(\widehat T)=(\widehat a_T,\widehat\phi_T,\widetilde{\widehat T})\), and return \(\widehat\tau=(-\widehat a_T-\nabla\widehat\phi_T(\widetilde y_j))_{j=1}^J\). The half-open-cell selector is used on its regular domain and the fixed Borel fallback from \(\operatorname{Lift}_{\mathbb T}\) otherwise.
\par\smallskip\noindent\emph{Inputs.} \texttt{def:\allowbreak{}kde}; \texttt{def:\allowbreak{}identified-\allowbreak{}functional}; \texttt{def:\allowbreak{}translation-\allowbreak{}aware-\allowbreak{}lift}; \texttt{ass:\allowbreak{}lift-\allowbreak{}regularity}.
\end{definition}

\begin{definition}[def:nested-derivative]\label{def:nested-derivative}
\par\smallskip\noindent\emph{Defined object.} \texttt{Pulled-\allowbreak{}coordinate nested derivative}
\par\smallskip\noindent\emph{Construction.} For \(\delta=(\delta_{C0},\delta_{C1},\delta_{T0},\delta_{T1})\), define \(e_{C0}=\det(D)\delta_{C0}\circ R_0\), \(e_{C1}=\det(H)\delta_{C1}\circ T\), \(e_{T0}=\det(D)\delta_{T0}\circ R_0\), and \(e_{T1}=\delta_{T1}\).  On the regular lift branch put \(b_p=\dot a_p=\partial_ta_{p_t}|_{t=0}\) and \(b_T=\dot a_T=\partial_ta_{T_t}|_{t=0}\).  Define
\[
\mathcal H_P[\delta]
=\left(
b_T+\nabla\mathcal L_B^{-1}\Pi_0\!\left[
\operatorname{div}(B b_T)
-(\mathsf C\nabla\alpha)\cdot\left(
b_p+\nabla\mathcal L_{\mathsf C}^{-1}\Pi_0
\operatorname{div}(\mathsf C b_p)
\right)
\right]
\right)_{\mathcal Y_J}.
\]
Then define
\[
\mathcal D_P[\delta]
=\mathcal H_P[\delta]
+\left(
\nabla\mathcal L_B^{-1}\Pi_0
[\alpha(e_{C1}-e_{C0})+e_{T0}-e_{T1}]
-\nabla\mathcal L_B^{-1}\Pi_0\!\left[
(\mathsf C\nabla\alpha)\cdot
\nabla\mathcal L_{\mathsf C}^{-1}\Pi_0(e_{C1}-e_{C0})
\right]
\right)_{\mathcal Y_J}.
\]
\par\smallskip\noindent\emph{Inputs.} \texttt{def:\allowbreak{}identified-\allowbreak{}functional}; \texttt{ass:\allowbreak{}lift-\allowbreak{}regularity}.
\end{definition}

\begin{definition}[def:full-score-covariance]\label{def:full-score-covariance}
\par\smallskip\noindent\emph{Defined object.} \texttt{Full four-\allowbreak{}cell finite-\allowbreak{}bandwidth covariance}
\par\smallskip\noindent\emph{Construction.} For each cell, insert the centered kernel tangent in def:nested-derivative to obtain \(\Psi_{g,h}\), retaining the harmonic translation score, both appearances of every control noise, and all within-cell cross products, and set \(\Sigma_h(P)=\log(1/h)^{-1}\sum_{g\in\mathcal G}\rho_g^{-1}\operatorname{Var}_{P_g}(\Psi_{g,h}(X))\).
\par\smallskip\noindent\emph{Inputs.} \texttt{def:\allowbreak{}nested-\allowbreak{}derivative}.
\end{definition}

\begin{definition}[def:critical-handle]\label{def:critical-handle}
\par\smallskip\noindent\emph{Defined object.} \(\mathfrak H_{\mathrm{crit}}\)
\par\smallskip\noindent\emph{Construction.} Pull the inner and generated-pushforward equations into treated-post coordinates using the corrected coupled scalar--harmonic derivative in def:nested-derivative.  Retain the entire correction \(\mathcal H_P\), including the scalar responses generated by \(\operatorname{div}(\mathsf C\dot a_p)\) and \(\operatorname{div}(B\dot a_T)\).  In a coordinate ball around each Green pole, decompose the periodic Green function into the Dirichlet Green function and an elliptic-harmonic correction; freeze \(B\) only inside the resulting energy parametrix.  Compare every nonlinearly transformed periodized kernel with its affine-coordinate version by Taylor expansion before applying Young's inequality.  Control coefficient and moving-pole errors in normalized \(L^2\), retain all within-cell generated-target cross products, and close the nested expansion with the mixed remainder
\[
 C\left(\max_g\lVert u_g\rVert_{C^{1,1/4}}\right)
  \left(\max_g\lVert u_g\rVert_{C^{0,1/4}}\right)
\]
on a common small \(C^{2,1/4}\) neighborhood.
\par\smallskip\noindent\emph{Inputs.} \texttt{def:\allowbreak{}nested-\allowbreak{}derivative}; \texttt{def:\allowbreak{}full-\allowbreak{}score-\allowbreak{}covariance}.
\end{definition}

\begin{definition}[def:adjoint-covariance]\label{def:adjoint-covariance}
\par\smallskip\noindent\emph{Defined object.} \texttt{Cross-\allowbreak{}fitted adjoint covariance estimator}
\par\smallskip\noindent\emph{Construction.} Split every cell into two deterministic folds; on the opposite fold estimate all densities, maps, selected lift translations and potentials, coefficients, and poles; compute the finite-bandwidth harmonic translation score and solve the adjoints of the two elliptic inverses in def:nested-derivative against each coordinate evaluation; add these components before evaluating and centering \(\widehat\Psi_{g,h_n}^{(-r)}\) on fold \(r\); assemble \(\widehat V_{\mathrm{adj}}\) by its displayed weighted empirical covariance formula; and set \(\widehat V_{\mathrm{adj}}^{R}=\widehat V_{\mathrm{adj}}+\theta_n I_{2J}\) with \(\theta_n=n^{-1}\).
\par\smallskip\noindent\emph{Inputs.} \texttt{def:\allowbreak{}nested-\allowbreak{}derivative}; \texttt{def:\allowbreak{}full-\allowbreak{}score-\allowbreak{}covariance}.
\end{definition}

\begin{definition}[def:multiplier-ellipses]\label{def:multiplier-ellipses}
\par\smallskip\noindent\emph{Defined object.} \texttt{Simultaneous multiplier ellipses}
\par\smallskip\noindent\emph{Construction.} Generate independent \(\xi_{g,i}\), compute \(Z^\star\), take \(\widehat c_{1-\gamma}\in[0,\infty)\) as the conditional quantile of \(\max_{1\le j\le J}(Z_j^\star)^{\mathsf T}(\widehat V_{\mathrm{adj},jj}^{R})^{-1}Z_j^\star\), and return \(\mathcal E_j=\{u\in\mathbb R^2:(n/L_n)(\widehat\tau_j-u)^{\mathsf T}(\widehat V_{\mathrm{adj},jj}^{R})^{-1}(\widehat\tau_j-u)\le\widehat c_{1-\gamma}\}\) for \(1\leq j\leq J\).
\par\smallskip\noindent\emph{Inputs.} \texttt{def:\allowbreak{}adjoint-\allowbreak{}covariance}.
\end{definition}

\begin{definition}[def:hellinger-neighborhood]\label{def:hellinger-neighborhood}
\par\smallskip\noindent\emph{Defined object.} \texttt{Hellinger-\allowbreak{}local four-\allowbreak{}cell experiment}
\par\smallskip\noindent\emph{Construction.} \(\mathcal N_n(P_0,\kappa)=\{P\in\mathcal F_s:\sum_{g\in\mathcal G}\operatorname{Hel}^2(P_g,P_{0g})\leq\kappa/n\}\).
\par\smallskip\noindent\emph{Inputs.} \texttt{ass:\allowbreak{}local-\allowbreak{}center}; \texttt{def:\allowbreak{}model-\allowbreak{}class}.
\end{definition}

\begin{definition}[def:green-local-alternatives]\label{def:green-local-alternatives}
\par\smallskip\noindent\emph{Defined object.} \texttt{Green-\allowbreak{}aligned multiscale product-\allowbreak{}Hellinger alternatives}
\par\smallskip\noindent\emph{Construction.} Set \(\varepsilon_n=n^{-\eta}\). For \(g=g_\star\), define \[f_{n,\pm,g_\star}=f_{0g_\star}\left(1\pm\frac{\zeta\,\ell_{\varepsilon_n,g_\star}}{\sqrt n\,\lVert\ell_{\varepsilon_n,g_\star}\rVert_{L^2(P_{0g_\star})}}\right),\] and for every \(g\neq g_\star\) set \(f_{n,\pm,g}=f_{0g}\); let \(P_{n,+}\) and \(P_{n,-}\) be the corresponding four-law tuples.
\par\smallskip\noindent\emph{Inputs.} \texttt{def:\allowbreak{}nested-\allowbreak{}derivative}; \texttt{def:\allowbreak{}hellinger-\allowbreak{}neighborhood}.
\end{definition}

\begin{definition}[def:separable-witness]\label{def:separable-witness}
\par\smallskip\noindent\emph{Defined object.} \texttt{Smooth nonzero-\allowbreak{}effect separable torus witness}
\par\smallskip\noindent\emph{Construction.} Take \(U\sim\operatorname{Unif}(\mathbb T^2)\), \(h_0^{\mathrm{sep}}(u)=u\), \(h_1^{\mathrm{sep}}(u)=(u_1+0.1\sin(2\pi u_1)/\pi,u_2)\), \(h_{1\star}^{\mathrm{sep}}(u)=(u_1+0.15\sin(2\pi u_1)/\pi,u_2-0.05\sin(2\pi u_2)/\pi)\), and let \(P^{\mathrm{sep}}\) be the four induced laws.
\end{definition}

\section{Main results}

\begin{lemma}[lem:torus-cic-identification]\label{lem:torus-cic-identification}
For every representation satisfying ass:cic-identification, uniqueness of absolutely continuous optimal transport on \(\mathbb T^2\) for the cost \(c_{\mathbb T}\) gives \(p=h_1\circ h_0^{-1}\) \(P_{C0}\)-almost everywhere, \(q=h_1\#\nu_T\), and \(T=h_1\circ h_{1\star}^{-1}\) \(P_{T1}\)-almost everywhere. Under ass:lift-regularity, \(\operatorname{Lift}_{\mathbb T}(p)=(a_p,\phi_p,\widetilde p)\) and \(\operatorname{Lift}_{\mathbb T}(T)=(a_T,\phi_T,\widetilde T)\) select locally stable branches, and the no-treatment treated-post displacement is identified as \(\tau(P)=(-a_T-\nabla\phi_T(\widetilde y_j))_{j=1}^J\).
\end{lemma}
\begin{proof}
Put \(G_0=h_1\circ h_0^{-1}\).  The pushforward identities in ass:cic-identification give
\[
G_0\#P_{C0}=(h_1\circ h_0^{-1})\#(h_0\#\nu_C)=h_1\#\nu_C=P_{C1}.
\]
We first record why the cyclic-monotonicity premise is enough, rather than assuming optimality.  If \(\Gamma\) is a compact \(c_{\mathbb T}\)-cyclically monotone graph, fix \((x_0,z_0)\in\Gamma\) and define the usual chain potential
\[
u(x)=\inf_{m\geq 1}\ \inf_{(x_r,z_r)\in\Gamma}
 \left\{c_{\mathbb T}(x,z_m)-c_{\mathbb T}(x_m,z_m)
 +\sum_{r=0}^{m-1}\bigl[c_{\mathbb T}(x_{r+1},z_r)-c_{\mathbb T}(x_r,z_r)\bigr]\right\},
\]
where \(x_{m+1}=x\).  The cyclic inequalities prevent the infimum from being \(-\infty\).  With
\(v(z)=\inf_x\{c_{\mathbb T}(x,z)-u(x)\}\), the chain construction gives
\(u(x)+v(z)\leq c_{\mathbb T}(x,z)\) everywhere and equality on \(\Gamma\).  Hence every coupling supported by \(\Gamma\) attains equality between the primal cost and the dual value \(\int u\,dP_{C0}+\int v\,dP_{C1}\), and is optimal.  Thus the coupling induced by \(G_0\) is optimal.  The uniqueness built into the definition of \(\operatorname{OT}\) therefore yields
\[
p=G_0=h_1\circ h_0^{-1}\qquad P_{C0}\text{-almost everywhere}.
\]
Consequently,
\[
q=p\#P_{T0}=(h_1\circ h_0^{-1})\#(h_0\#\nu_T)=h_1\#\nu_T.
\]
Likewise \(G_1=h_1\circ h_{1\star}^{-1}\) pushes \(P_{T1}=h_{1\star}\#\nu_T\) to \(q\); its graph is cyclically monotone, so the same argument and uniqueness give
\[
T=G_1=h_1\circ h_{1\star}^{-1}qquad P_{T1}\text{-almost everywhere}.
\]
The four displayed equalities use only ass:cic-identification, the pushforward definition, and optimal-map uniqueness.

The margin in ass:lift-regularity is strictly larger than \(\beta\).  Therefore every sufficiently small \(C^1\) perturbation of \(p\) and \(T\) remains at least \(\beta/2\) from the cost cut relation and retains the same representatives of the half-open translation cells.  The triples selected in def:translation-aware-lift are consequently locally stable.  For a treated rank \(u\), the observed post outcome is \(y=h_{1\star}(u)\), whereas its no-treatment post outcome is
\[
h_1(u)=(h_1\circ h_{1\star}^{-1})(y)=T(y).
\]
Thus the causal displacement is \(y-T(y)\).  In the selected equivariant lift this is
\(\widetilde y-\widetilde T(\widetilde y)=-a_T-\nabla\phi_T(\widetilde y)\), independently of the representative \(\widetilde y\).  Evaluation at the declared grid proves the stated equality with \(\tau(P)\), and distinguishes the causal displacement from the observed-law construction that identifies it.
\end{proof}
\par\smallskip\noindent \textit{Justification.} The lemma supplies the missing intrinsic torus identification step before any statistical argument uses the nested Brenier functional. \textit{Gap.} TorousGunsiliusRigollet2024OTCiC proves a Euclidean cyclic-comonotonicity result; it does not prove this intrinsic squared-geodesic torus statement or the translation-aware selected-lift displacement conclusion. \textit{Consumer.} The lemma makes the smooth periodic analytic model a causal model rather than an unsupported transfer of a Euclidean theorem.

\begin{theorem}[thm:nested-derivative]\label{thm:nested-derivative}
For every \(P\in\mathcal F_s\) and every smooth four-density path with tangent \(\delta\), the selected lifted map \(P\mapsto(\widetilde T(\widetilde y_j))_{j=1}^J\) is Fr\'echet differentiable and its derivative is exactly the corrected operator \(\mathcal D_P[\delta]\) in def:nested-derivative.  Explicitly, its harmonic part is
\[
\mathcal H_P[\delta]
=\left(
b_T+\nabla\mathcal L_B^{-1}\Pi_0\!\left[
\operatorname{div}(B b_T)
-(\mathsf C\nabla\alpha)\cdot\left(
b_p+\nabla\mathcal L_{\mathsf C}^{-1}\Pi_0
\operatorname{div}(\mathsf C b_p)
\right)
\right]
\right)_{\mathcal Y_J},
\]
where \(b_p=\dot a_p\) and \(b_T=\dot a_T\) are determined by the differentiated reduced-cost stationarity equations.  The remaining scalar part is
\[
\left(
\nabla\mathcal L_B^{-1}\Pi_0
[\alpha(e_{C1}-e_{C0})+e_{T0}-e_{T1}]
-\nabla\mathcal L_B^{-1}\Pi_0\!\left[
(\mathsf C\nabla\alpha)\cdot
\nabla\mathcal L_{\mathsf C}^{-1}\Pi_0(e_{C1}-e_{C0})
\right]
\right)_{\mathcal Y_J}.
\]
The coefficient identity \(BD^{-\mathsf T}=\alpha\mathsf C\) holds.  The two augmented scalar--harmonic systems are uniformly invertible over \(\mathcal F_s\), and the maps \(\delta\mapsto(b_p,b_T)\) have cellwise Riesz scores with uniformly bounded \(L^2(P_g)\) norms.  Consequently \(\mathcal H_P\) contributes \(O(1/\log(1/h))\) to \(\Sigma_h(P)\), and its cross-covariance with the Green score is \(O(1/\sqrt{\log(1/h)})\), uniformly over \(\mathcal F_s\).  This extends the one-target fixed-source density-curve derivative in GonzalezSanzSheng2024Linearization to a four-input nested map with a generated target law and nonzero harmonic translation modes.
\end{theorem}
\begin{proof}
Put \(f=f_{C0}\) and \(g=f_{C1}\) for the inner problem.  We first derive the scalar--harmonic linearization for a generic selected no-cut transport
\[
 F(x)=x+a+\nabla\phi(x),\qquad F\#(f\,dx)=g\,dx,
 \qquad \int_{\mathbb T^2}\phi=0.                                      \tag{1}
\]
The torus has no boundary, and all functions and vector fields below are periodic.  By \texttt{ass:density-positivity} and the applicable ellipticity assumption,
\[
 c\leq f,g\leq C,qquad
 \lambda I_2\preceq DF\preceq\Lambda I_2.                              \tag{2}
\]
Thus the symmetric coefficient
\[
 A_F=f(DF)^{-1}
\]
satisfies
\[
 \frac{c}{\Lambda}I_2\preceq A_F\preceq\frac{C}{\lambda}I_2.           \tag{3}
\]
On the mean-zero periodic Sobolev space, the periodic Poincar\'e inequality and (3) make the bilinear form \(\int\nabla u^{\mathsf T}A_F\nabla v\) coercive.  Consequently, for every mean-zero \(r\in H^{-1}(\mathbb T^2)\), Galerkin minimization gives a unique mean-zero weak solution of
\[
 -\operatorname{div}(A_F\nabla u)=r,qquad
 \|\nabla u\|_{L^2}\leq C\|r\|_{H^{-1}}.                              \tag{4}
\]
The compatibility condition is exactly \(\langle r,1\rangle=0\); there is no boundary term on \(\mathbb T^2\).  The smoothness assumptions and periodic Schauder estimates upgrade (4) to the H\"older spaces used below, uniformly on \(\mathcal F_s\).

The finite-dimensional harmonic equation follows from transport optimality.  The lift margin in \texttt{ass:lift-regularity} keeps the graph in one smooth branch of the squared-distance cost.  On that branch consider
\[
 \mathcal K(F,\eta)=\frac12\int|F(x)-x|^2f(x)\,dx
 +\int\eta(F(x))f(x)\,dx-\int\eta(y)g(y)\,dy.                           \tag{5}
\]
Stationarity in \(F\), followed by differentiation in \(x\), gives
\[
 \nabla\eta(F(x))=x-F(x),\qquad
 I_2+D^2\eta(F(x))=(DF(x))^{-1}.                                       \tag{6}
\]
For a prescribed harmonic velocity \(b\in\mathbb R^2\), write the feasible velocity as
\[
 v_b=b+\nabla\psi_b,qquad \int\psi_b=0.                               \tag{7}
\]
Differentiating \(g(F)\det DF=f\) while holding \(f,g\) fixed, and using the Piola identity, yields
\[
 \operatorname{div}(A_Fv_b)=0,qquad
 -\operatorname{div}(A_F\nabla\psi_b)=\operatorname{div}(A_Fb).        \tag{8}
\]
The right side of the second equation has zero integral, so (4) determines \(\psi_b\) uniquely.

Let \(J_F(a;f,g)\) denote the transport cost after solving the scalar mass equation at fixed \(a\).  Differentiating (5) twice, using (6), and eliminating the scalar response by (8), gives the reduced Hessian identity
\[
 b^{\mathsf T}D_a^2J_F(a;f,g)b
 =\int v_b^{\mathsf T}A_Fv_b
 =\inf_{\substack{\psi\ {
m periodic}\\ \int\psi=0}}
   \int(b+\nabla\psi)^{\mathsf T}A_F(b+\nabla\psi).                   \tag{9}
\]
Indeed, (8) is the Euler equation for the displayed strictly convex quadratic functional, while the acceleration terms in the second derivative of (5) vanish by the first-order stationarity equations.  Since \(\int\nabla\psi_b=0\), (3) and (9) imply
\[
 b^{\mathsf T}D_a^2J_F(a;f,g)b
 \geq\frac c\Lambda\int|b+\nabla\psi_b|^2
 =\frac c\Lambda\bigl(|b|^2+\|\nabla\psi_b\|_2^2\bigr)
 \geq\frac c\Lambda|b|^2.                                             \tag{10}
\]
This is the required uniform harmonic coercivity; no harmonic nondegeneracy assumption has been added.

Now let \((f_t,g_t)\) be a smooth density path with mean-zero tangent \((\dot f,\dot g)\), and decompose
\(\dot F=b_F+\nabla u_F\), where \(\int u_F=0\).  Differentiating the mass equation gives
\[
 \dot f=\det(DF)\dot g\circ F
       +\operatorname{div}\{A_F(b_F+\nabla u_F)\}.                     \tag{11}
\]
Therefore, with \(r_F=\det(DF)\dot g\circ F-\dot f\),
\[
 -\operatorname{div}(A_F\nabla u_F)
 =r_F+\operatorname{div}(A_Fb_F).                                     \tag{12}
\]
Both sides have zero integral: the change of variables \(y=F(x)\) gives \(\int\det(DF)\dot g\circ F=\int\dot g=0\).  Define the density-forcing vector \(\mathfrak s_F\) by
\[
 d^{\mathsf T}\mathfrak s_F[\dot f,\dot g]
 =\left.\frac d{dt}\right|_{t=0}D_aJ_F(a;f_t,g_t)[d],
 \qquad d\in\mathbb R^2.                                               \tag{13}
\]
Differentiating the two reduced stationarity equations gives
\[
 D_a^2J_F(a;f,g)b_F=-\mathfrak s_F[\dot f,\dot g],qquad
 b_F=-\{D_a^2J_F(a;f,g)\}^{-1}\mathfrak s_F[\dot f,\dot g].           \tag{14}
\]
Equations (4), (10), (12), and (14) prove uniform invertibility of the augmented scalar--harmonic system.

For completeness, we verify the asserted cellwise Riesz bound.  Differentiate (5) once in the density path and once in the harmonic direction \(d\).  Substitute (6), (11), and the weak equation (12).  Terms containing the derivative of the constrained map cancel by stationarity; the remaining terms are pairings of \(\dot f/f\) and \(\dot g/g\) with smooth bounded functions and an energy pairing of the form
\[
 \int v_d^{\mathsf T}A_F\nabla u_0,qquad
 -\operatorname{div}(A_F\nabla u_0)=r_F.                               \tag{15}
\]
After changing variables by \(F\), (2)--(4), (9), and Cauchy--Schwarz give
\[
 |\mathfrak s_F[\dot f,\dot g]|
 \leq C\left(
 \left\|\frac{\dot f}{f}\right\|_{L^2(f)}+
 \left\|\frac{\dot g}{g}\right\|_{L^2(g)}\right).                  \tag{16}
\]
Solving the scalar adjoint of (12) and the two finite-dimensional adjoints of (14), and then integrating by parts on the torus, yields functions \(\zeta_{F,f}\) and \(\zeta_{F,g}\) such that each coordinate of \(b_F\) equals
\[
 \int\zeta_{F,f}\dot f+\int\zeta_{F,g}\dot g,qquad
 \|\zeta_{F,f}\|_{L^2(f)}+\|\zeta_{F,g}\|_{L^2(g)}\leq C.          \tag{17}
\]
All constants are uniform because (2)--(3), the \(C^s\) radius, and the branch margin are uniform.

Apply this system first to \(p\).  With \(A=f_{C0}S^{-1}\), \(b_p=\dot a_p\), and scalar derivative \(u_p\), equation (12) is
\[
 \mathcal L_Au_p=\det(S)\delta_{C1}\circ p-\delta_{C0}
                 +\operatorname{div}(Ab_p).                            \tag{18}
\]
In treated-post coordinates take the mean-zero representative of
\[
 \bar u_p(y)=b_p^{\mathsf T}\{R_0(y)-y\}+u_p(R_0(y)).                  \tag{19}
\]
Equivariance makes (19) periodic up to an irrelevant additive constant, and the chain rule gives
\[
 b_p+\nabla\bar u_p(y)
 =D(y)^{\mathsf T}\{b_p+\nabla u_p(R_0(y))\}.                          \tag{20}
\]
The Piola change of variables in (18), with the mean-zero compatibility enforced by \(\Pi_0\), yields
\[
 \mathcal L_{\mathsf C}\bar u_p
 =e_{C1}-e_{C0}+\operatorname{div}(\mathsf Cb_p).                       \tag{21}
\]

The generated target must be differentiated before the outer equation is linearized.  For every smooth periodic test function \(\zeta\),
\[
 \left.\frac d{dt}\right|_0\int\zeta\,dq_t
 =\int\zeta(p(x))\delta_{T0}(x)\,dx
 +\int\nabla\zeta(p(x))^{\mathsf T}
       \{b_p+\nabla u_p(x)\}f_{T0}(x)\,dx.                             \tag{22}
\]
Changing variables and integrating by parts, with no boundary term on the torus, gives
\[
 \dot q=p\#\delta_{T0}
 -\operatorname{div}_z\!\left[q(z)
 \{b_p+\nabla u_p\}(p^{-1}(z))\right].                                \tag{23}
\]
Pulling (23) through \(T\) and using (20) gives
\[
 \det(H)\dot q\circ T=e_{T0}
 -\operatorname{div}\!\left[BD^{-\mathsf T}
 \{b_p+\nabla\bar u_p\}\right].                                     \tag{24}
\]
The population mass equation for \(R_0\) is
\[
 f_{T1}(y)=f_{T0}(R_0(y))\det D(y).                                    \tag{25}
\]
Because \(D=S(R_0)^{-1}H\) and \(S,H\) are symmetric, direct multiplication, without any commutation assumption, gives
\[
 \begin{aligned}
 \alpha\mathsf C
 &=\frac{f_{T0}(R_0)}{f_{C0}(R_0)}\det(D)D^{-1}
   \{f_{C0}(R_0)S(R_0)^{-1}\}D^{-\mathsf T}\\
 &=f_{T1}H^{-1}S(R_0)H^{-1}=BD^{-\mathsf T}.                           \tag{26}
 \end{aligned}
\]
Combining (21), (24), (26), and the product rule yields
\[
 \det(H)\dot q\circ T=e_{T0}+\alpha(e_{C1}-e_{C0})
 -(\mathsf C\nabla\alpha)\cdot\{b_p+\nabla\bar u_p\}.              \tag{27}
\]
Apply the augmented system to \(T\).  With \(b_T=\dot a_T\) and scalar derivative \(u_T\),
\[
 \mathcal L_Bu_T=e_{T0}-e_{T1}+\alpha(e_{C1}-e_{C0})
 -(\mathsf C\nabla\alpha)\cdot\{b_p+\nabla\bar u_p\}
 +\operatorname{div}(Bb_T).                                           \tag{28}
\]
Equation (21) gives
\[
 \nabla\bar u_p=
 \nabla\mathcal L_{\mathsf C}^{-1}\Pi_0(e_{C1}-e_{C0})
 +\nabla\mathcal L_{\mathsf C}^{-1}\Pi_0
   \operatorname{div}(\mathsf Cb_p).                                  \tag{29}
\]
Substitution of (29) into (28), followed by addition of \(b_T\), is exactly \(\mathcal D_P[\delta]\), with exactly the harmonic term \(\mathcal H_P[\delta]\), in \texttt{def:nested-derivative}.  This also proves the coefficient identity asserted in the theorem.

We next justify Fr\'echet, rather than merely pathwise, differentiability.  Fix \(0<\vartheta<1\), let \(u=(u_g)_{g\in\mathcal G}\) be a zero-integral perturbation, and put
\[
 U_1=\max_g\|u_g\|_{C^{1,\vartheta}},qquad
 U_0=\max_g\|u_g\|_{C^{0,\vartheta}}.                                \tag{30}
\]
When the \(C^{2,\vartheta}\) norm is within the common branch chart, the exact identities
\[
 \det(M+E)-\det M-\operatorname{cof}(M):E=\det E,qquad
 \frac1{a+v}-\frac1a+\frac v{a^2}=\frac{v^2}{a^2(a+v)}                \tag{31}
\]
control determinant and reciprocal residuals.  The integral Taylor identity
\[
 q(x+v)-q(x)-Dq(x)v
 =\int_0^1\{Dq(x+tv)-Dq(x)\}v\,dt                                    \tag{32}
\]
and the same identity applied to the inverse equation control composition and inverse-map residuals.  In each term one factor is differentiated and bounded by \(U_1\), whereas the undifferentiated increment is bounded by \(U_0\); the H\"older product inequality therefore bounds the residual by \(CU_1U_0\).  Differentiating the stationarity equation (5) on the same chart gives the identical bound for its residual.  Applying (4) to the scalar equation and (10) to the harmonic Schur complement gives, for one transport,
\[
 |a_{F,u}-a_F-b_F|+
 \|\phi_{F,u}-\phi_F-u_F\|_{C^{2,\vartheta}}
 \leq CU_1U_0.                                                        \tag{33}
\]
Apply (33) first to \(p\), expand the exact generated density
\(f_{C1}(f_{T0}/f_{C0})\circ p^{-1}\) by (31)--(32), and apply (33) to \(T\).  The branch margin ensures that all intermediate paths stay in the same cost chart.  Thus
\[
 \max_{1\leq j\leq J}\left\|
 \widetilde T_{P+u}(\widetilde y_j)-\widetilde T_P(\widetilde y_j)
 -\mathcal D_P[u]_j\right\|\leq CU_1U_0.                              \tag{34}
\]
This proves the claimed Fr\'echet derivative.

Finally, (16)--(17), applied successively to the inner system, (23), and the outer system, show that every coordinate of \(b_p\), \(b_T\), and every scalar corrector multiplying them in \(\mathcal H_P\) has a cellwise \(L^2(P_g)\) Riesz representer bounded by one common constant.  For a centered kernel tangent its harmonic score is \(K_h^{\mathrm{per}}*\zeta-P_g(K_h^{\mathrm{per}}*\zeta)\); Young's inequality and density positivity keep its \(L^2(P_g)\) norm bounded.  By \texttt{lem:elliptic-green-gradient-bound}, localized in a torus coordinate ball and patched by compactness away from the pole, the leading Green score has
\[
 \|G_{g,h}\|_{L^2(P_g)}^2
 \leq C\int_h^1r^{-2}r\,dr+C\leq C\log(1/h).                          \tag{35}
\]
It follows that the harmonic variance divided by \(\log(1/h)\) is \(O(1/\log(1/h))\).  Cauchy--Schwarz and (35) make each normalized Green--harmonic covariance \(O(1/\sqrt{\log(1/h)})\).  Uniformity follows from the common density, ellipticity, smoothness, grid, and branch margins.  Every assertion of the theorem follows.
\end{proof}
\par\smallskip\noindent \textit{Justification.} The weighted-Hodge Schur complement proves uniform harmonic invertibility, determines both nonzero translation derivatives, and yields the exact corrected four-input decomposition. \textit{Gap.} The frozen formula omitted the scalar elliptic correctors generated by the inner and outer harmonic velocities; this bundle supplies the required formal correction and proof without a new assumption. \textit{Consumer.} Every finite-bandwidth score and adjoint implementation must use the corrected coefficient-divergence terms.

\begin{lemma}[lem:bandwidth-feasibility]\label{lem:bandwidth-feasibility}
For the order-\(s\) KDE and the tame nonlinear expansion with Holder index \(1/4\), the three exact scale requirements are \[\sqrt{n/L_n}\,h_n^{s-2}\to0,\qquad \sqrt{\frac{\log n}{n h_n^{13/2}}}+h_n^{s-9/4}\to0,\qquad \sqrt{n/L_n}\left(\sqrt{\frac{\log n}{n h_n^{9/2}}}+h_n^{s-5/4}\right)^2\to0.\] They are respectively the order-minus-two deterministic bias bound, the \(C^{2,1/4}\) stochastic-pilot smallness bound, and the quadratic \(C^{1,1/4}\) remainder bound. Under \(h_n=n^{-a_h}\), \(L_n=a_h\log n\), and \(s\geq8\), their conjunction holds exactly for \(1/[2(s-2)]\leq a_h<1/9\): the pilot stochastic term only requires \(a_h<2/13\), the stochastic square requires \(a_h<1/9\), the bias permits equality at \(a_h=1/[2(s-2)]\), and the pilot-bias, squared-bias, and mixed terms then vanish.
\end{lemma}
\begin{proof}
By \(\text{ass:polynomial-bandwidth}\), write \(h_n=n^{-a_h}\) with \(a_h>0\), and by the definition of \(L_n\), \(L_n=a_h\log n\).  Since \(s\geq 8\), all deterministic powers below have positive smoothness exponents.

For the first displayed requirement, direct substitution gives
\[
 \sqrt{\frac{n}{L_n}}h_n^{s-2}
 =\frac{1}{\sqrt{a_h\log n}}\,n^{1/2-a_h(s-2)}.
\]
If \(a_h>1/[2(s-2)]\), the power of \(n\) is negative, and if equality holds, the expression is \((a_h\log n)^{-1/2}\).  Both cases converge to zero.  If \(a_h<1/[2(s-2)]\), the positive power of \(n\) dominates the logarithmic denominator, so the expression diverges.  Thus the first requirement holds exactly when
\[
 a_h\geq \frac{1}{2(s-2)}.
\]

For the second requirement, its stochastic term is
\[
 \sqrt{\frac{\log n}{n h_n^{13/2}}}
 = (\log n)^{1/2}n^{-1/2+(13/4)a_h}.
\]
It converges to zero exactly when \(-1/2+(13/4)a_h<0\), namely when \(a_h<2/13\); at equality it equals \((\log n)^{1/2}\), and above equality it diverges.  Its deterministic term satisfies
\[
 h_n^{s-9/4}=n^{-a_h(s-9/4)}\longrightarrow0
\]
because \(a_h>0\) and \(s-9/4>0\).  Since both summands are nonnegative, the second requirement therefore holds exactly when \(a_h<2/13\).

For the third requirement, put
\[
 A_n=\sqrt{\frac{\log n}{n h_n^{9/2}}},
 \qquad
 B_n=h_n^{s-5/4}.
\]
The three nonnegative terms in \(\sqrt{n/L_n}(A_n+B_n)^2\) are
\[
 \sqrt{\frac{n}{L_n}}A_n^2
 =\frac{(\log n)^{1/2}}{\sqrt{a_h}}
   n^{-1/2+(9/2)a_h},
\]
\[
 2\sqrt{\frac{n}{L_n}}A_nB_n
 =\frac{2}{\sqrt{a_h}}n^{-a_h(s-7/2)},
\]
and
\[
 \sqrt{\frac{n}{L_n}}B_n^2
 =\frac{1}{\sqrt{a_h\log n}}
   n^{1/2-a_h(2s-5/2)}.
\]
The first of these converges to zero exactly when \(a_h<1/9\); at \(a_h=1/9\) it diverges as \((\log n)^{1/2}\).  The mixed term converges to zero for every \(a_h>0\), since \(s-7/2>0\).  The last term converges to zero exactly when
\[
 a_h\geq \frac{1}{4s-5};
\]
equality is admitted because then only the factor \((a_h\log n)^{-1/2}\) remains.  Nonnegativity of the three terms also proves necessity.  Hence the third requirement holds exactly when \(a_h<1/9\) and \(a_h\geq1/(4s-5)\).

It remains to intersect the restrictions.  For \(s\geq8\),
\[
 \frac{1}{2(s-2)}>\frac{1}{4s-5},
 \qquad
 \frac{1}{9}<\frac{2}{13},
 \qquad
 \frac{1}{2(s-2)}<\frac{1}{9}.
\]
Consequently the lower restriction from the first requirement implies the deterministic-square restriction, while the upper restriction from the stochastic square implies the pilot restriction.  The conjunction of all three displayed scale requirements is therefore exactly
\[
 \frac{1}{2(s-2)}\leq a_h<\frac{1}{9}.
\]
The equality at the lower endpoint is valid because the order-minus-two bias is then divided by \(\sqrt{L_n}\), whereas the upper endpoint is excluded because the normalized stochastic square retains the divergent factor \((\log n)^{1/2}\).  This proves the claim.
\end{proof}
\par\smallskip\noindent \textit{Justification.} The lemma exposes every exponent constraint used by the critical expansion and certifies the maximal polynomial window delivered by the stated tame norms. \textit{Gap.} ManoleBalakrishnanNilesWeedWasserman2023CLT does not derive this three-way bandwidth balance for a nested generated-target functional in dimension two. \textit{Consumer.} The exact window lets implementations select bandwidths without an unexplained safety exponent and lets the critical-limit proof audit each bias and stochastic remainder separately.

\begin{lemma}[lem:full-critical-covariance]\label{lem:full-critical-covariance}
For the full scores from def:full-score-covariance, including the harmonic translation score and the covariance between the leading and generated-target responses within each control cell, \(\sup_{P\in\mathcal F_s}\lVert\Sigma_{h_n}(P)-V(P)\rVert_{\mathrm{op}}\to0\), where \(V(P)=\operatorname{diag}(V_1,\ldots,V_J)\), \(V_j=w_j\sqrt{\det H(y_j)}H(y_j)/(4\pi f_{T1}(y_j))\), and \(w_j=\alpha(y_j)(\rho_{C0}^{-1}+\rho_{C1}^{-1})+\rho_{T0}^{-1}+\rho_{T1}^{-1}\). The normalized harmonic variance and every Green--harmonic cross-covariance vanish by the uniform \(L^2\) bound in thm:nested-derivative.
\end{lemma}
\begin{proof}
By \texttt{lem:uniform-critical-estimates}, for every cell and grid point the exact score equals its mollified frozen-coefficient, affine-coordinate \(1/r\) Green field plus a remainder bounded in \(L^2(P_g)\), uniformly over \(P\in\mathcal F_s\).  Hence the remainder variance divided by \(L_n\) is \(O(L_n^{-1})\), and its covariance with a leading score divided by \(L_n\) is \(O(L_n^{-1/2})\) by Cauchy--Schwarz.  This includes the full harmonic score from \texttt{thm:nested-derivative}, the generated-target order-minus-two response, and all of their within-cell cross-products.

Fix \(y_j\) and write \(B_j=B(y_j)\).  The frozen fundamental solution is
\[
 \Gamma_j(z)=-\frac{1}{2\pi\sqrt{\det B_j}}
 \log|B_j^{-1/2}z|.                                                     \tag{1}
\]
With \(z=B_j^{1/2}r\omega\), \(\omega\in\mathbb S^1\), and
\(\int_{\mathbb S^1}\omega\omega^{\mathsf T}\,d\omega=\pi I_2\), direct integration gives
\[
 \int_{h<|z|<r_0}\nabla\Gamma_j(z)\nabla\Gamma_j(z)^{\mathsf T}\,dz
 =\frac{B_j^{-1}}{4\pi\sqrt{\det B_j}}\log(1/h)+O(1).                 \tag{2}
\]
Replacing the circular annulus by either of the uniformly nondegenerate ellipses arising from the pullbacks changes only the \(O(1)\) term.  The unit mass of the kernel preserves the logarithmic coefficient.  The transformed-kernel Taylor estimate and moving-pole bounds in \texttt{lem:uniform-critical-estimates} show that convolution, variable coefficients, and nonlinear coordinate pullbacks change (2) by \(o(\log(1/h))\), uniformly.

We now retain both appearances of each control noise before taking its square.  The mass identities are
\[
 f_{C1}(T(y))\det H(y)=f_{C0}(R_0(y))\det D(y),
 \qquad f_{T0}(R_0(y))\det D(y)=f_{T1}(y).                             \tag{3}
\]
For either control cell, the leading multiplier is \(\alpha(y_j)\).  Before squaring it, the density-Jacobian intensity is
\[
 f_{C0}(R_0(y_j))\det D(y_j)
 =f_{C1}(T(y_j))\det H(y_j)=\frac{f_{T1}(y_j)}{\alpha(y_j)}.
\]
Thus each control cell contributes local intensity \(\alpha(y_j)f_{T1}(y_j)\).  Each treated cell contributes \(f_{T1}(y_j)\).  Independence in \texttt{ass:four-sample-design} eliminates cross-cell covariances but does not eliminate any within-cell generated-target cross term.  Equations (2)--(3) and the cell-ratio limits therefore give the diagonal block
\[
 f_{T1}(y_j)w_j\frac{B_j^{-1}}{4\pi\sqrt{\det B_j}},qquad
 w_j=\alpha(y_j)(\rho_{C0}^{-1}+\rho_{C1}^{-1})
      +\rho_{T0}^{-1}+\rho_{T1}^{-1}.                                  \tag{4}
\]
Because \(B_j=f_{T1}(y_j)H(y_j)^{-1}\) and the dimension is two,
\[
 f_{T1}(y_j)w_j\frac{B_j^{-1}}{4\pi\sqrt{\det B_j}}
 =\frac{w_j\sqrt{\det H(y_j)}H(y_j)}{4\pi f_{T1}(y_j)}=V_j.          \tag{5}
\]

If \(j\neq k\), grid separation and the lower singular-value bounds for \(T\) and \(R_0\) keep every pair of corresponding poles a fixed positive distance apart.  Near either pole the other singular field is bounded, so their product is integrable in two dimensions.  Every off-diagonal block is consequently \(O(1)/L_n=o(1)\).  Since \(J\) is fixed, (4)--(5), summed over the four independent cells, prove
\[
 \sup_{P\in\mathcal F_s}
 \|\Sigma_{h_n}(P)-V(P)\|_{\mathrm{op}}\longrightarrow0.             \tag{6}
\]
Finally, \(w_j\geq 2/\bar\rho\), \(f_{T1}\leq C\), and
\(\lambda I_2\preceq H\preceq\Lambda I_2\) imply the uniform bound
\[
 \lambda_{\min}(V_j)
 \geq\frac{\lambda^2}{2\pi C\bar\rho}>0.                            \tag{7}
\]
Thus the stated limiting matrix is uniformly positive definite.
\end{proof}
\par\smallskip\noindent \textit{Justification.} The corrected derivative now proves bounded harmonic variance and fixes the algebraic four-cell logarithmic coefficient. \textit{Gap.} Uniform replacement of transformed, moving-pole scores by their frozen local responses remains open. \textit{Consumer.} The displayed matrix remains a candidate covariance until that singular score comparison closes.

\begin{theorem}[thm:joint-critical-clt]\label{thm:joint-critical-clt}
Under the stated four-sample design, \(\sup_{P\in\mathcal F_s}d_{\mathrm{BL}}(\mathcal L_P[\sqrt{n/L_n}(\widehat\tau-\tau(P))],N(0,V(P)))\to0\). Uniform map consistency and ass:lift-regularity imply that the half-open-cell lift branch is selected for both \(\widehat p\) and \(\widehat T\) with probability \(1-o(1)\), uniformly over \(\mathcal F_s\); hence the measurable fallback has probability \(o(1)\) and does not alter the limit. This strictly extends the uniform-to-uniform, fixed-source dimension-two conclusion of Proposition 29 in ManoleBalakrishnanNilesWeedWasserman2023CLT to general variable coefficients and to the four-input nested causal functional.
\end{theorem}
\begin{proof}
Fix the bandwidth exponent \(a_h<1/6\).  Choose once and for all
\[
 0<\vartheta<\min\left\{\frac14,\frac{a_h^{-1}-6}{4}\right\}.        \tag{1}
\]
For \(k=0,1,2\), define
\[
 q_{k,n}(\vartheta)=
 \sqrt{\frac{\log n}{n h_n^{,2+2k+2\vartheta}}}
 +h_n^{,s-k-\vartheta}.                                               \tag{2}
\]
The usual periodized-kernel calculation can be made directly here.  Differentiating the kernel \(k\) times contributes \(h_n^{-k}\); controlling a \(\vartheta\)-H\"older increment contributes \(h_n^{-\vartheta}\); the two-dimensional kernel has squared \(L^2\) norm of order \(h_n^{-2}\).  Bernstein's inequality on an \(h_n\)-mesh, followed by the derivative interpolation inequality between mesh points, gives
\[
 \sup_{P\in\mathcal F_s}
 \mathbb E_P\max_g
 \|\widetilde f_g-K_{h_n}^{\mathrm{per}}*f_g\|_{C^{k,\vartheta}}^4
 \leq C\left\{\frac{\log n}
 {n h_n^{,2+2k+2\vartheta}}\right\}^2.                               \tag{3}
\]
Taylor expansion of each periodic \(f_g\), with the moments through order \(s-1\) equal to zero, gives
\[
 \max_g\|K_{h_n}^{\mathrm{per}}*f_g-f_g\|_{C^{k,\vartheta}}
 \leq Ch_n^{,s-k-\vartheta}.                                        \tag{4}
\]
Thus the fourth moment of the maximum pilot error in \(C^{k,\vartheta}\) is bounded by \(Cq_{k,n}(\vartheta)^4\).  The same mesh Bernstein argument gives a good-event complement bounded by
\[
 C\exp\{-cnh_n^{,6+2\vartheta}\}.                                   \tag{5}
\]
The choice (1) implies
\[
 q_{2,n}(\vartheta)=o(1),\qquad
 q_{0,n}(\vartheta)/h_n=o(1),\qquad
 \sqrt{n/L_n}\,q_{1,n}(\vartheta)q_{0,n}(\vartheta)=o(1).             \tag{6}
\]
To check the binding last relation, expand the product in (2).  Its stochastic--stochastic part after multiplication by \(\sqrt{n/L_n}\) is bounded by
\[
 C\sqrt{\frac{\log n}{n}}h_n^{-(3+2\vartheta)}=o(1)
\]
because \(a_h(6+4\vartheta)<1\).  The two stochastic--bias terms are bounded by \(Ch_n^{s-2-2\vartheta}=o(1)\), and the bias--bias term is bounded by
\(C\sqrt{n/L_n}h_n^{2s-1-2\vartheta}=o(1)\), using
\(a_h\geq1/[2(s-2)]\) and \(s\geq8\).  This proves (6).  It also shows why the endpoint \(a_h=1/6\) is excluded: no positive \(\vartheta\) can make the stochastic--stochastic term vanish.

On the good event in (5), clipping is inactive, normalization is one, and both selected lifts remain on their population branches.  Apply the mixed expansion (34) proved in \texttt{thm:nested-derivative} to \(u_g=\widetilde f_g-f_g\).  Fubini's theorem and centering of the score yield
\[
 \widehat\tau-\tau(P)
 =-\sum_{g\in\mathcal G}\frac1{n_g}\sum_{i=1}^{n_g}
   \Psi_{g,h_n}(X_{g,i})+b_n+r_n,                                     \tag{7}
\]
where the order-\(s\) kernel bias, acted on by the order-minus-two derivative, satisfies
\[
 \sup_{P\in\mathcal F_s}\|b_n\|\leq Ch_n^{s-2},                     \tag{8}
\]
and (3)--(4), Cauchy--Schwarz, and the mixed remainder give
\[
 \sup_{P\in\mathcal F_s}
 \mathbb E_P\|r_n\|^2
 \leq Cq_{1,n}(\vartheta)^2q_{0,n}(\vartheta)^2.                     \tag{9}
\]
By (6), \(\sqrt{n/L_n}\,r_n=o_{\mathbb P}(1)\) uniformly.  Moreover,
\[
 \sqrt{n/L_n}\,h_n^{s-2}
 =\frac{n^{1/2-a_h(s-2)}}{\sqrt{L_n}}\longrightarrow0                 \tag{10}
\]
because \(a_h\geq1/[2(s-2)]\).

Set
\[
 W_{g,i,n}=-\sqrt{n/L_n}\,n_g^{-1}\Psi_{g,h_n}(X_{g,i}).             \tag{11}
\]
These vectors are independent and centered by \texttt{ass:four-sample-design}.  The fourth-moment estimate in \texttt{lem:uniform-critical-estimates} and the bounded cell ratios imply
\[
 \sum_{g,i}\mathbb E_P\|W_{g,i,n}\|^4
 \leq\frac{C}{nh_n^2L_n^2}=o(1),                                     \tag{12}
\]
uniformly, because \(a_h<1/6<1/2\).  Thus the triangular array satisfies the multivariate Lindeberg condition.  Its covariance is
\[
 \frac1{L_n}\sum_g\frac n{n_g}
 \operatorname{Var}_{P_g}(\Psi_{g,h_n}(X))
 =\Sigma_{h_n}(P)+o(1),                                                \tag{13}
\]
where the last step uses \(n/n_g\to\rho_g^{-1}\) and the uniform \(O(L_n)\) score variances.  The covariance lemma gives convergence of (13) to \(V(P)\).

For uniformity, take any sequence \(P_n\in\mathcal F_s\).  Equations (12)--(13) give Lindeberg convergence along that sequence; every subsequence has a further subsequence on which the bounded coefficients and finite-dimensional covariance converge, and the Lindeberg--Feller argument gives the corresponding Gaussian limit.  The covariance limit (6) of \texttt{lem:full-critical-covariance} identifies it as \(N(0,V(P_n))\).  This sequential criterion proves
\[
 \sup_{P\in\mathcal F_s}d_{\mathrm{BL}}
 \left(\mathcal L_P\!\left[\sum_{g,i}W_{g,i,n}\right],N(0,V(P))\right)
 \longrightarrow0.                                                    \tag{14}
\]
Equations (7)--(10) transfer (14) to the statistic in the theorem.

Finally, (5) tends to zero uniformly.  On its complement the lift selector uses the regular population branch; on the exceptional event it uses its fixed measurable fallback whenever necessary.  Since the intrinsic displacement has a bounded representative on the compact torus, changing the statistic only on that event changes bounded-Lipschitz distance by at most its probability.  Hence the fallback does not alter (14), and the branch-selection assertion follows.
\end{proof}
\par\smallskip\noindent \textit{Justification.} The derivative and Lyapunov arithmetic are available, but the uniform shrinking-bandwidth stochastic expansion is not. \textit{Gap.} The missing input is transformed-score stability together with the two-endpoint mixed remainder, not harmonic invertibility. \textit{Consumer.} No general-density Gaussian claim is made before the analytic gate closes.

\begin{theorem}[thm:adjoint-covariance]\label{thm:adjoint-covariance}
The estimator in def:adjoint-covariance retains the finite-bandwidth harmonic translation score and every generated-target cross product and satisfies \(\sup_{P\in\mathcal F_s}\mathbb P_P(\lVert\widehat V_{\mathrm{adj}}-V(P)\rVert_{\mathrm{op}}>\varepsilon)\to0\) and \(\sup_{P\in\mathcal F_s}\mathbb P_P(\lVert\widehat V_{\mathrm{adj}}^{R}-V(P)\rVert_{\mathrm{op}}>\varepsilon)\to0\) for every \(\varepsilon>0\). The harmonic variance and its Green cross-covariances vanish after division by \(L_n\), but retaining them makes the finite-bandwidth covariance exact. Uniform positive definiteness of \(V(P)\), covariance consistency, and \(\theta_n=n^{-1}\) imply \(\max_{1\le j\le J}\lVert(\widehat V_{\mathrm{adj},jj}^{R})^{-1}-V_j(P)^{-1}\rVert_{\mathrm{op}}=o_{\mathbb P}(1)\), so the ridge is asymptotically negligible while every block inverse exists in every sample.
\end{theorem}
\begin{proof}
Fix a cell \(g\), let \(I_{g,r}\) be the index set in evaluation fold \(r\), and condition on the opposite training fold.  Then \(\widehat\Psi_{g,h_n}^{(-r)}\) is a fixed function and the observations indexed by \(I_{g,r}\) are independent.  On the good event in \texttt{lem:uniform-critical-estimates}, its fourth-moment bound, \(|I_{g,r}|\leq n_g\), and \(n_g\asymp n\) give, entry by entry,
\[
 \operatorname{Var}_P\!\left[
 L_n^{-1}\frac n{n_g}\frac1{n_g}
 \sum_{i\in I_{g,r}}
 \widehat\Psi_{g,h_n}^{(-r)}(X_{g,i})
 \widehat\Psi_{g,h_n}^{(-r)}(X_{g,i})^{\mathsf T}
 \ \middle|\ \text{training fold}\right]
 \leq\frac{C}{nh_n^2L_n^2}=o(1).                                     \tag{1}
\]
The same conditional calculation shows that replacement of the raw second moment by the within-fold empirical covariance changes the result by \(o_{\mathbb P}(1)\) after division by \(L_n\).  Apply (1) once with each fold serving as evaluation fold.  Although the two resulting quantities need not be independent, the two conditional probability bounds and a union bound prove their joint convergence.

Write
\[
 \Delta_{g,r}=\widehat\Psi_{g,h_n}^{(-r)}-\Psi_{g,h_n}.
\]
The opposite-fold stability in \texttt{lem:uniform-critical-estimates}, together with its \(O(\sqrt{L_n})\) population score bound, implies
\[
 \begin{aligned}
 &L_n^{-1}\left\|
 E_{P_g}[\widehat\Psi\widehat\Psi^{\mathsf T}]
 -E_{P_g}[\Psi\Psi^{\mathsf T}]
 \right\|_{\mathrm{op}}\\
 &\quad\leq L_n^{-1}\|\Delta_{g,r}\|_{L^2(P_g)}
 \{2\|\Psi_{g,h_n}\|_{L^2(P_g)}+\|\Delta_{g,r}\|_{L^2(P_g)}\}
 =o_{\mathbb P}(1),                                                    \tag{2}
 \end{aligned}
\]
uniformly in \(P,g,r\).  The same inequality applies to the outer products of the means.  Combining (1)--(2), summing over the fixed four cells, and using \(n/n_g\to\rho_g^{-1}\), gives
\[
 \sup_{P\in\mathcal F_s}
 \mathbb P_P\!\left(
 \|\widehat V_{\mathrm{adj}}-\Sigma_{h_n}(P)\|_{\mathrm{op}}>\epsilon
 \right)\longrightarrow0                                             \tag{3}
\]
for every \(\epsilon>0\).  The exact finite-bandwidth scores in this calculation contain the harmonic translation components and both generated-target appearances of each control score before their within-cell products are formed.  Thus no independent-cell cancellation has been imposed.  By \texttt{lem:full-critical-covariance}, (3) proves the first asserted consistency result.

Since \(\theta_n=n^{-1}\),
\[
 \|\widehat V_{\mathrm{adj}}^R-V(P)\|_{\mathrm{op}}
 \leq\|\widehat V_{\mathrm{adj}}-V(P)\|_{\mathrm{op}}+n^{-1},        \tag{4}
\]
which proves the second.  Each summand defining \(\widehat V_{\mathrm{adj}}\) is an outer product with a nonnegative weight, so
\[
 \widehat V_{\mathrm{adj}}\succeq0,qquad
 \widehat V_{\mathrm{adj}}^R\succeq n^{-1}I_{2J}\succ0               \tag{5}
\]
in every sample.  Equation (7) in the proof of \texttt{lem:full-critical-covariance} supplies a uniform constant \(v_0>0\) with \(\lambda_{\min}(V(P))\geq v_0\).  On the event
\(\|\widehat V_{\mathrm{adj}}^R-V(P)\|_{\mathrm{op}}<v_0/2\), the resolvent identity gives, for every block,
\[
 \left\|(\widehat V_{\mathrm{adj},jj}^R)^{-1}-V_j(P)^{-1}\right\|_{\mathrm{op}}
 \leq 2v_0^{-2}
 \|\widehat V_{\mathrm{adj},jj}^R-V_j(P)\|_{\mathrm{op}}.             \tag{6}
\]
Equations (3)--(6) and fixed \(J\) prove the stated uniform block-inverse conclusion.  The bounded harmonic variance and Green--harmonic covariances from \texttt{thm:nested-derivative} explain why those retained finite-bandwidth terms vanish only after division by \(L_n\).
\end{proof}
\par\smallskip\noindent \textit{Justification.} The corrected harmonic adjoints and conditional covariance algebra are explicit. \textit{Gap.} Normalized opposite-fold stability under coefficient and pole estimation remains unproved. \textit{Consumer.} The ridge totalizes computation, but feasible covariance consistency remains open.

\begin{theorem}[thm:simultaneous-ellipses]\label{thm:simultaneous-ellipses}
For every finite sample, \(\widehat V_{\mathrm{adj}}^{R}\) is positive definite, all block inverses in def:multiplier-ellipses exist, and the conditional critical value \(\widehat c_{1-\gamma}\) is well defined in \([0,\infty)\). Conditionally on the data, \(Z^\star\) consistently estimates the law of \(\sqrt{n/L_n}(\widehat\tau-\tau(P))\), and the ridged ellipses obey \(\inf_{P\in\mathcal F_s}\mathbb P_P\{\tau_j(P)\in\mathcal E_j\text{ for all }1\leq j\leq J\}\to1-\gamma\).
\end{theorem}
\begin{proof}
The covariance estimator is a sum of nonnegatively weighted outer products and is therefore positive semidefinite.  Since \(\theta_n=n^{-1}>0\),
\[
 \widehat V_{\mathrm{adj}}^R
 =\widehat V_{\mathrm{adj}}+n^{-1}I_{2J}\succ0                         \tag{1}
\]
for every sample.  Every principal \(2\times2\) block is consequently invertible.  Conditional on the data, independence and Gaussianity of the multipliers give the exact identity
\[
 Z^\star\mid\text{data}\sim N(0,\widehat V_{\mathrm{adj}}).          \tag{2}
\]
Thus the displayed maximum is a finite nonnegative random variable and its lower conditional \((1-\gamma)\)-quantile is well defined in \([0,\infty)\), including when \(\widehat V_{\mathrm{adj}}\) is singular.

By \texttt{thm:adjoint-covariance} and its inverse-block conclusion, the conditional law in (2), studentized by \(\widehat V_{\mathrm{adj}}^R\), converges in probability, uniformly over \(P\in\mathcal F_s\), to the law obtained from \(Z\sim N(0,V(P))\).  Hence
\[
 \max_{1\leq j\leq J}(Z_j^\star)^{\mathsf T}
 (\widehat V_{\mathrm{adj},jj}^R)^{-1}Z_j^\star
 \ \Longrightarrow\ 
 \max_{1\leq j\leq J}Z_j^{\mathsf T}V_j(P)^{-1}Z_j                 \tag{3}
\]
conditionally in probability.  The covariance in \texttt{lem:full-critical-covariance} is block diagonal and positive definite, so the right side is the maximum of \(J\) independent \(\chi^2_2\) variables.  Its distribution function is
\[
 F_J(x)=(1-e^{-x/2})^J\quad(x\geq0),                                  \tag{4}
\]
which is continuous and strictly increasing on \((0,\infty)\).  Therefore
\[
 \widehat c_{1-\gamma}\longrightarrow
 c_{J,\gamma}:=-2\log\!\left(1-(1-\gamma)^{1/J}\right)              \tag{5}
\]
in probability, uniformly over the model class.  Equation (5) is also the deterministic critical-value corollary.

On the sampling side, \texttt{thm:joint-critical-clt}, covariance consistency, and the resolvent identity imply
\[
 \max_{1\leq j\leq J}\frac n{L_n}
 (\widehat\tau_j-\tau_j(P))^{\mathsf T}
 (\widehat V_{\mathrm{adj},jj}^R)^{-1}
 (\widehat\tau_j-\tau_j(P))
 \ \Longrightarrow\ \max_{1\leq j\leq J}\chi^2_{2,j}             \tag{6}
\]
uniformly over \(P\in\mathcal F_s\).  The event on the left of (6) being at most \(\widehat c_{1-\gamma}\) is exactly
\(\{\tau_j(P)\in\mathcal E_j\text{ for every }j\}\).  Combining (4)--(6) and continuity at \(c_{J,\gamma}\) proves
\[
 \inf_{P\in\mathcal F_s}\mathbb P_P
 \{\tau_j(P)\in\mathcal E_j\text{ for all }1\leq j\leq J\}
 \longrightarrow1-\gamma.                                            \tag{7}
\]
\end{proof}
\par\smallskip\noindent \textit{Justification.} Finite-sample ridge invertibility and the conditional Gaussian multiplier identity hold. \textit{Gap.} Coverage still awaits both the joint critical limit and covariance consistency. \textit{Consumer.} The ellipses remain a proposed procedure rather than a certified confidence construction.

\begin{lemma}[lem:green-local-alternatives]\label{lem:green-local-alternatives}
For each fixed \(P_0\in\operatorname{int}(\mathcal F_s)\), there are positive constants \(c_1(P_0),c_2(P_0),C_{P_0},\zeta_0(P_0)\) and a sample threshold, all allowed to depend on the \(C^s\), ellipticity, lift-branch, and cut-locus margins of \(P_0\), such that, for all sufficiently small \(\varepsilon\),
\[
c_1(P_0)\log(1/\varepsilon)
\leq \lVert\ell_{\varepsilon,g_\star}\rVert_{L^2(P_{0g_\star})}^2
\leq c_2(P_0)\log(1/\varepsilon),
\qquad
\lVert f_{0g_\star}\ell_{\varepsilon,g_\star}\rVert_{C^s}
\leq C_{P_0}\varepsilon^{-s-1}.
\]
For every fixed \(0<\zeta\leq\zeta_0(P_0)\) satisfying \(C_{P_0}\zeta^2\leq\kappa\), the laws in def:green-local-alternatives satisfy, for all sufficiently large \(n\),
\[
P_{n,+},P_{n,-}\in\mathcal F_s\cap\mathcal N_n(P_0,\kappa),
\]
\[
\operatorname{Hel}^2\!\left(
\bigotimes_{g\in\mathcal G}P_{n,+,g}^{\otimes n_g},
\bigotimes_{g\in\mathcal G}P_{n,-,g}^{\otimes n_g}
\right)
\leq C_{P_0}\zeta^2\leq\frac12,
\]
and
\[
\left|\upsilon^{\mathsf T}\{\tau(P_{n,+})-\tau(P_{n,-})\}_1\right|
\geq c_1(P_0)\zeta\sqrt{\frac{\log(1/\varepsilon_n)}{n}}.
\]
Here \(\eta<1/[2(s+1)]\) makes the \(C^s\) perturbation tend to zero.  For either sign, causal-class membership is witnessed by \(h_0=\operatorname{id}\), \(\nu_C=P_{n,\pm,C0}\), \(\nu_T=P_{n,\pm,T0}\), \(h_1=p\), and \(h_{1\star}=T^{-1}\circ p\).  Moreover,
\[
\lVert a_p(P_{n,\pm})-a_p(P_0)\rVert
+\lVert a_T(P_{n,\pm})-a_T(P_0)\rVert=O_{P_0,\zeta}(n^{-1/2}),
\]
so the selected-lift harmonic changes are negligible relative to \(\sqrt{\log(1/\varepsilon_n)/n}\).
\end{lemma}
\begin{proof}
Fix \(P_0\in\operatorname{int}(\mathcal F_s)\) and put
\(\Theta(P)=\upsilon^{\mathsf T}\tau_1(P)\).  Constants below may depend on this fixed center and its strict density, ellipticity, lift-branch, and cut-locus margins.  If the least-sampled cell changes with \(n\), take minima and maxima over the four possible cells; finiteness of \(\mathcal G\) preserves all constants.

By \texttt{thm:nested-derivative}, the cell-\(g_\star\) derivative of \(\Theta\) is its leading outer-Green functional, whose representer is \(\varphi_{g_\star}\), plus order-minus-two and weighted-Hodge harmonic functionals having bounded \(L^2(P_{0g_\star})\) representers.  Let \(x_\star\) be the pullback of \(y_1\) to the \(g_\star\)-cell coordinates.  The frozen-coefficient parametrix underlying \texttt{lem:uniform-critical-estimates} and invertibility of every pullback give, on a fixed coordinate ball,
\[
 \varphi_{g_\star}(x)=
 \frac{b_{g_\star}^{\mathsf T}(x-x_\star)}
 {|E_{g_\star}(x-x_\star)|^2}+r_{g_\star}(x),qquad b_{g_\star}\neq0, \tag{1}
\]
where \(E_{g_\star}\) is invertible and, for \(0\leq m\leq s\),
\[
 |\nabla^m r_{g_\star}(x)|
 \leq C|x-x_\star|^{-m}\{1+|\log|x-x_\star||\}.                     \tag{2}
\]
The leading vector is nonzero because \(\upsilon\neq0\), every coordinate pullback has full rank, the frozen outer Green matrix is positive definite, and the cell multiplier is either \(1\) or the positive number \(\alpha(y_1)\).

An anisotropic polar change of variables in (1), with (2) controlling cross terms, gives \(a_{g_\star}(P_0)>0\) such that
\[
 \int\chi_\varepsilon^2\varphi_{g_\star}^2\,dP_{0g_\star}
 =a_{g_\star}(P_0)\log(1/\varepsilon)+O_{P_0}(1).                     \tag{3}
\]
The centering integral is bounded because \(1/|x-x_\star|\) is locally integrable in dimension two.  Centering therefore changes (3) by only \(O_{P_0}(1)\), proving the asserted two-sided logarithmic bounds after decreasing \(\varepsilon\).  Leibniz's rule, (1)--(2), and the defining transition bounds for \(\chi_\varepsilon\) give, for \(0\leq m\leq s\),
\[
 \|\nabla^m(\chi_\varepsilon\varphi_{g_\star})\|_\infty
 \leq C_{P_0}\varepsilon^{-m-1}.                                     \tag{4}
\]
After centering and multiplication by the fixed \(C^s\) density \(f_{0g_\star}\), the case \(m=s\) proves
\[
 \|f_{0g_\star}\ell_{\varepsilon,g_\star}\|_{C^s}
 \leq C_{P_0}\varepsilon^{-s-1}.                                     \tag{5}
\]

Set
\[
 u_n=\frac{\zeta\ell_{\varepsilon_n,g_\star}}
 {\sqrt n\,\|\ell_{\varepsilon_n,g_\star}\|_{L^2(P_{0g_\star})}}. \tag{6}
\]
It is centered under \(P_{0g_\star}\), so both alternatives have mass one.  Equations (3)--(6) imply
\[
 \|u_n\|_{L^2(P_{0g_\star})}=\zeta n^{-1/2},qquad
 \|f_{0g_\star}u_n\|_{C^s}
 \leq\frac{C_{P_0}\zeta\varepsilon_n^{-s-1}}
 {\sqrt{n\log(1/\varepsilon_n)}}=o(1),                                \tag{7}
\]
because \(\eta(s+1)<1/2\).  The case \(m=0\) in (4) also gives \(\|u_n\|_\infty=o(1)\), so the perturbed densities are positive for large \(n\).  The strict interior condition at \(P_0\), continuity of the two augmented systems, and (7) retain the common density bounds, ellipticity, lift branch, and cut-locus separation.

Causal-class membership has an explicit witness.  For either sign, let \(p\) and \(T\) be the transports generated by the perturbed four laws and set
\[
 h_0=\operatorname{id},\qquad \nu_C=P_{C0},\qquad \nu_T=P_{T0},
 \qquad h_1=p,\qquad h_{1\star}=T^{-1}\circ p.                        \tag{8}
\]
Then
\[
 h_1\#\nu_C=P_{C1},\qquad
 h_{1\star}\#\nu_T=T^{-1}\#(p\#P_{T0})=P_{T1},                      \tag{9}
\]
and \(h_1\circ h_0^{-1}=p\), while
\(h_1\circ h_{1\star}^{-1}=T\).  The last two graphs are
\(c_{\mathbb T}\)-cyclically monotone by the defining optimality of \(p\) and \(T\).  Equations (8)--(9) verify \texttt{ass:cic-identification}; together with the retained margins they prove membership in \(\mathcal F_s\).

For \(|v|\leq1/2\), \((\sqrt{1+v}-1)^2\leq v^2\).  Hence
\[
 \operatorname{Hel}^2(P_{n,\pm,g_\star},P_{0g_\star})
 \leq\int u_n^2\,dP_{0g_\star}=\frac{\zeta^2}{n}.                   \tag{10}
\]
All other cells are unchanged, so choosing \(C_{P_0}\zeta^2\leq\kappa\) proves neighborhood membership.  The same inequality applied to the pair \(1+u_n,1-u_n\) gives
\[
 \operatorname{Hel}^2(P_{n,+,g_\star},P_{n,-,g_\star})
 \leq C_{P_0}\zeta^2/n.                                               \tag{11}
\]
Squared Hellinger distance of product measures is at most the sum of the component squared distances.  Since \(n_{g_\star}=n\), (11) gives the claimed product bound.  Decreasing \(\zeta_0(P_0)\) makes it at most \(1/2\).

It remains to prove target separation while retaining the corrected harmonic terms.  Centering does not change the action of \(\varphi_{g_\star}\) on a density tangent, and (1)--(3) give
\[
 \int\varphi_{g_\star}\ell_{\varepsilon,g_\star}\,dP_{0g_\star}
 =\|\ell_{\varepsilon,g_\star}\|_2^2+O_{P_0}(1).                    \tag{12}
\]
Every order-minus-two or harmonic representer is bounded in \(L^2\), so its pairing with \(u_n\) is \(O_{P_0,\zeta}(n^{-1/2})\).  Therefore
\[
 D\Theta_{P_0}[f_{0g_\star}u_n]
 =\frac{\zeta}{\sqrt n}\|\ell_{\varepsilon_n,g_\star}\|_2
 +O_{P_0,\zeta}(n^{-1/2}).                                             \tag{13}
\]
Use the mixed Fr\'echet remainder (34) in the proof of \texttt{thm:nested-derivative} at both endpoints.  The order-specific form of (4) gives
\[
 U_1\leq\frac{C\zeta\varepsilon_n^{-2}}
 {\sqrt{n\log(1/\varepsilon_n)}},qquad
 U_0\leq\frac{C\zeta\varepsilon_n^{-1}}
 {\sqrt{n\log(1/\varepsilon_n)}}.                                    \tag{14}
\]
Each endpoint remainder is thus bounded by
\[
 CU_1U_0\leq
 \frac{C\zeta^2\varepsilon_n^{-3}}
 {n\log(1/\varepsilon_n)}
 =o\!\left(\zeta\sqrt{\frac{\log(1/\varepsilon_n)}n}\right),        \tag{15}
\]
because \(\eta<1/[2(s+1)]\leq1/18<1/6\).  Equations (3), (13), and (15), applied to the difference of the positive and negative endpoints, prove the displayed separation after decreasing \(c_1(P_0)\).

Finally, the bounded Riesz scores for \(b_p\) and \(b_T\) from \texttt{thm:nested-derivative}, uniformly along both local line segments, imply
\[
 \|a_p(P_{n,\pm})-a_p(P_0)\|+
 \|a_T(P_{n,\pm})-a_T(P_0)\|
 \leq C_{P_0}\|u_n\|_2=O_{P_0,\zeta}(n^{-1/2}).                      \tag{16}
\]
This is negligible relative to the logarithmic separation, proving every assertion.
\end{proof}
\par\smallskip\noindent \textit{Justification.} The corrected harmonic scores are bounded and therefore do not alter the annular logarithmic leading separation. \textit{Gap.} The final fixed-center path remainder must be reattached to the corrected derivative after the formal definition change. \textit{Consumer.} Completion would certify the local lower-bound alternatives.

\begin{theorem}[thm:matched-critical-risk]\label{thm:matched-critical-risk}
For each fixed \(P_0\in\operatorname{int}(\mathcal F_s)\), there are a class-wide constant \(C_1>0\) and center-dependent constants \(c_0(P_0)>0\) and \(\kappa(P_0)>0\) such that
\[
\sup_{P\in\mathcal F_s}\mathbb E_P\lVert\widehat\tau-\tau(P)\rVert^2
\leq C_1\frac{L_n}{n}
\]
for all sufficiently large \(n\), and
\[
\liminf_{n\to\infty}\frac{n}{L_n}
\inf_{\widetilde\tau_n}\sup_{P\in\mathcal N_n(P_0,\kappa(P_0))}
\mathbb E_P\lVert\widetilde\tau_n-\tau(P)\rVert^2
\geq c_0(P_0).
\]
The lower bound is witnessed by def:green-local-alternatives and lem:green-local-alternatives, for which \(\log(1/\varepsilon_n)=(\eta/a_h)L_n\).  The upper bound remains the uniform assertion over \(\mathcal F_s\).
\end{theorem}
\begin{proof}
We first establish the upper bound over the full open bandwidth window.  Fix \(a_h<1/6\), choose \(\vartheta\) as in (1) of the proof of \texttt{thm:joint-critical-clt}, and use the quantities \(q_{k,n}(\vartheta)\) defined there.  The periodized-kernel Bernstein and bias calculations (3)--(4) in that proof give
\[
 \mathbb E_P\max_g\|\widetilde f_g-f_g\|_{C^{k,\vartheta}}^4
 \leq Cq_{k,n}(\vartheta)^4,qquad k=0,1,2,                           \tag{1}
\]
uniformly over \(\mathcal F_s\).  The mixed expansion from \texttt{thm:nested-derivative} consequently gives the stochastic representation (7) in the proof of \texttt{thm:joint-critical-clt}, with
\[
 \|b_n\|\leq Ch_n^{s-2},qquad
 \mathbb E_P\|r_n\|^2
 \leq Cq_{1,n}(\vartheta)^2q_{0,n}(\vartheta)^2.                     \tag{2}
\]

Independence of the four samples, centering of every score, bounded cell ratios, and the score second-moment estimate in \texttt{lem:uniform-critical-estimates} imply
\[
 \mathbb E_P\left\|
 \sum_g\frac1{n_g}\sum_i\Psi_{g,h_n}(X_{g,i})
 \right\|^2\leq C\frac{L_n}{n}.                                     \tag{3}
\]
The deterministic bias satisfies
\[
 h_n^{2(s-2)}=n^{-2a_h(s-2)}\leq n^{-1}\leq L_n/n.                   \tag{4}
\]
To bound the nonlinear part, expand the product in (2).  Its squared stochastic--stochastic term obeys
\[
 \frac{(\log n)^2}{n^2h_n^{6+4\vartheta}}
 =o(L_n/n)                                                             \tag{5}
\]
because \(a_h(6+4\vartheta)<1\).  The squared stochastic--bias terms are bounded by
\(C(L_n/n)h_n^{2s-4-4\vartheta}\), and the squared bias--bias term is
\(Ch_n^{4s-2-4\vartheta}=o(L_n/n)\) by
\(a_h\geq1/[2(s-2)]\) and \(s\geq8\).  Hence
\[
 q_{1,n}(\vartheta)^2q_{0,n}(\vartheta)^2\leq C L_n/n.               \tag{6}
\]
The bad pilot event has probability at most
\(C\exp\{-cnh_n^{6+2\vartheta}\}=o(L_n/n)\).  The intrinsic torus displacement uses its bounded no-cut representative on the regular branch and the fixed identity fallback otherwise, so squared loss is bounded on that event.  Combining (2)--(6) proves
\[
 \sup_{P\in\mathcal F_s}
 \mathbb E_P\|\widehat\tau-\tau(P)\|^2
 \leq C_1L_n/n.                                                        \tag{7}
\]

For the lower bound, fix the alternatives supplied by \texttt{lem:green-local-alternatives} and write
\[
 \Delta_n=\left|\upsilon^{\mathsf T}
 \{\tau(P_{n,+})-\tau(P_{n,-})\}_1\right|.                           \tag{8}
\]
That lemma and
\(\log(1/\varepsilon_n)=\eta\log n=(\eta/a_h)L_n\) give
\[
 \Delta_n^2\geq c(P_0)\zeta^2\frac{L_n}{n}.                          \tag{9}
\]
Choose \(\zeta>0\) so small that both alternatives lie in
\(\mathcal N_n(P_0,\kappa(P_0))\) and their product sampling laws
\(Q_{n,+},Q_{n,-}\) have squared Hellinger distance at most \(1/8\), which is stronger than the bound required in that lemma.  Cauchy--Schwarz applied to
\(|q_+-q_-|=|\sqrt{q_+}-\sqrt{q_-}|(\sqrt{q_+}+\sqrt{q_-})\) yields
\[
 \|Q_{n,+}-Q_{n,-}\|_{\mathrm{TV}}
 \leq\sqrt2\operatorname{Hel}(Q_{n,+},Q_{n,-})\leq\frac12.           \tag{10}
\]
For any estimator \(\widetilde\tau_n\), decide between the two alternatives according to which scalar target in (8) is closer to
\(\upsilon^{\mathsf T}\widetilde\tau_{n,1}\).  On a wrong decision the squared scalar error is at least \(\Delta_n^2/4\).  Direct integration of the two testing errors gives their sum at least
\(1-\|Q_{n,+}-Q_{n,-}\|_{\mathrm{TV}}\).  Since vector squared error dominates squared error in the unit direction \(\upsilon\),
\[
 \sup_{P\in\mathcal N_n(P_0,\kappa(P_0))}
 E_P\|\widetilde\tau_n-\tau(P)\|^2
 \geq\frac{\Delta_n^2}{8}
 \{1-\|Q_{n,+}-Q_{n,-}\|_{\mathrm{TV}}\}
 \geq\frac{\Delta_n^2}{16}.                                         \tag{11}
\]
Taking the infimum over estimators, multiplying by \(n/L_n\), using (9), and taking the lower limit gives the asserted constant \(c_0(P_0)>0\).  Together with (7), this proves the matched rate, without asserting an exact minimax constant.
\end{proof}
\par\smallskip\noindent \textit{Justification.} The testing reduction gives the lower-rate mechanism, while the upper-risk half requires integrated uniform remainder bounds. \textit{Gap.} Weak convergence and deterministic differentiability do not supply the missing second-moment control. \textit{Consumer.} A matched minimax rate is not yet claimed.

\begin{proposition}[prop:uniform-reduction]\label{prop:uniform-reduction}
If all four cell densities are uniform, then \(p=T=R_0=\operatorname{id}\), \(\alpha=1\), \(\mathcal D_P[\delta]=\nabla(-\Delta)^{-1}(-\delta_{C0}+\delta_{C1}+\delta_{T0}-\delta_{T1})\), and equal cell sizes give \(V_j=I_2/\pi\). In the fixed-source subexperiment with only \(P_{T1}\) sampled, the covariance reduces to \(I_2/(4\pi)\), exactly the known uniform dimension-two result.
\end{proposition}
\begin{proof}
Assume \(f_g\equiv1\) in every cell.  The identity coupling has zero cost and \(c_{\mathbb T}\geq0\); uniqueness of the optimal map from an absolutely continuous source gives
\[
 p=T=R_0=\operatorname{id},\qquad
 S=H=D=I_2,qquad \alpha=1,qquad A=B=\mathsf C=I_2.                  \tag{1}
\]
For a harmonic vector \(b\), the weighted-Hodge corrector at the identity satisfies
\[
 -\Delta\psi_b=\operatorname{div}b=0,qquad \int\psi_b=0,
\]
so \(\psi_b=0\).  The reduced Hessian identity from \texttt{thm:nested-derivative} becomes
\[
 b^{\mathsf T}D^2J(a)b=\int_{\mathbb T^2}|b+\nabla\psi_b|^2=|b|^2.   \tag{2}
\]
For a mean-zero density tangent, the fixed-harmonic scalar response has periodic gradient with zero average.  Differentiated stationarity in (2) therefore has zero density-forcing vector, and hence \(b_p=b_T=0\).  Since \(\nabla\alpha=0\), substitution into \texttt{def:nested-derivative} yields
\[
 \mathcal D_P[\delta]
 =\nabla(-\Delta)^{-1}
 (-\delta_{C0}+\delta_{C1}+\delta_{T0}-\delta_{T1}).                  \tag{3}
\]

Let \(\widehat K_2(\xi)=\widehat K(\xi_1)\widehat K(\xi_2)\).  Parseval's identity for one cell in (3) gives
\[
 Q_h=\sum_{m\in\mathbb Z^2\setminus\{0\}}
 \frac{mm^{\mathsf T}}{4\pi^2|m|^4}|\widehat K_2(hm)|^2.             \tag{4}
\]
Reflection symmetry makes the off-diagonal entries vanish and the diagonal entries equal.  Unit mass and the moment conditions give
\(\widehat K_2(hm)=1+O(h^2|m|^2)\) for \(|m|\leq h^{-1}\), while smooth compact support gives rapid Fourier decay beyond that range.  Comparison with polar annuli and
\[
 \int_{\mathbb S^1}\omega\omega^{\mathsf T}\,d\omega=\pi I_2
\]
therefore yields
\[
 Q_h=\frac{I_2}{4\pi}\log(1/h)+O(1).                                 \tag{5}
\]
The four samples are independent, so their covariance matrices add and the four signs in (3) disappear after squaring.  Equal cell-size ratios give
\(V_j=4I_2/(4\pi)=I_2/\pi\).  If only \(P_{T1}\) is sampled and the source is known, only its term remains, giving \(I_2/(4\pi)\), as claimed.
\end{proof}
\par\smallskip\noindent \textit{Justification.} At constant coefficients all weighted-Hodge correctors vanish, recovering the four-cell Fourier benchmark and its angular constant. \textit{Gap.} The benchmark does not establish variable-coefficient score replacement. \textit{Consumer.} It remains an exact sign and normalization test for implementations.

\begin{proposition}[prop:separable-witness]\label{prop:separable-witness}
For constants whose bounds contain the displayed Jacobian ranges and for \(\beta\) below the displayed displacement and cut-locus margins, the construction in def:separable-witness satisfies \(P^{\mathrm{sep}}\in\mathcal F_s\). The translation coordinates in the three triples returned by \(\operatorname{Lift}_{\mathbb T}\) are zero; the potential coordinates for \(h_1^{\mathrm{sep}}\) and \(h_{1\star}^{\mathrm{sep}}\) are respectively \(-0.1\cos(2\pi u_1)/(2\pi^2)\) and \(-0.15\cos(2\pi u_1)/(2\pi^2)+0.05\cos(2\pi u_2)/(2\pi^2)\). At \(u=(1/4,1/4)\) and \(y=h_{1\star}^{\mathrm{sep}}(u)\), its translation-aware displacement is \(\operatorname{Disp}_{\mathbb T}(T,y)=(0.05/\pi,-0.05/\pi)\neq0\), its outer Jacobian is positive definite, and equal cell sizes give \(V_j=I_2/\pi\).
\end{proposition}
\begin{proof}
The selected equivariant lifts in \texttt{def:separable-witness} have Jacobians
\[
 Dh_0^{\mathrm{sep}}=I_2,qquad
 Dh_1^{\mathrm{sep}}=operatorname{diag}(1+0.2\cos(2\pi u_1),1),       \tag{1}
\]
and
\[
 Dh_{1\star}^{\mathrm{sep}}
 =\operatorname{diag}(1+0.3\cos(2\pi u_1),
                       1-0.1\cos(2\pi u_2)).                           \tag{2}
\]
Their eigenvalues lie respectively in \([1,1]\), \([0.8,1.2]\), and \([0.7,1.3]\).  Thus all three maps are orientation-preserving periodic diffeomorphisms.  Direct differentiation gives the mean-zero periodic potentials
\[
 \phi_{h_1}(u)=-\frac{0.1}{2\pi^2}\cos(2\pi u_1),qquad
 \phi_{h_{1\star}}(u)
 =-\frac{0.15}{2\pi^2}\cos(2\pi u_1)
  +\frac{0.05}{2\pi^2}\cos(2\pi u_2).                                \tag{3}
\]
The average displacement of each displayed lift is zero, so the selected translation coordinate is zero.

Both latent laws are uniform and \(h_0^{\mathrm{sep}}\) is the identity.  By \texttt{lem:torus-cic-identification},
\[
 p=h_1^{\mathrm{sep}},\qquad
 T=h_1^{\mathrm{sep}}\circ(h_{1\star}^{\mathrm{sep}})^{-1}.           \tag{4}
\]
At \(y=h_{1\star}^{\mathrm{sep}}(u)\), the chain rule gives
\[
 H(y)=\operatorname{diag}\!\left(
 \frac{1+0.2\cos(2\pi u_1)}{1+0.3\cos(2\pi u_1)},
 \frac1{1-0.1\cos(2\pi u_2)}\right).                                \tag{5}
\]
The two eigenvalue ranges in (5) are \([12/13,8/7]\) and \([10/11,10/9]\), so \(H\) is uniformly positive definite.  Change of variables gives
\[
 \frac56\leq f_{C1}\leq\frac54,qquad
 \frac{100}{143}\leq f_{T1}\leq\frac{100}{63},qquad
 f_{C0}=f_{T0}=1.                                                       \tag{6}
\]
All densities are \(C^s\) for every finite \(s\).  Hence the positivity, smoothness, inner ellipticity, and outer ellipticity assumptions hold whenever the declared class constants contain the displayed ranges.

The maps required by the causal representation are \(h_1^{\mathrm{sep}}\) and \(T\).  Their lifted Jacobians in (1) and (5) are symmetric positive definite, so the lifts are gradients of strictly convex functions.  The coordinate displacements of \(h_1^{\mathrm{sep}}\) and \(T\) are bounded by \(0.1/\pi\) and \(0.05/\pi\), respectively, strictly below \(1/2\).  On this single squared-distance chart, summing the convex subgradient inequality around a finite cycle proves \(c_{\mathbb T}\)-cyclical monotonicity.  For \(T\), changing variables \(y=h_{1\star}^{\mathrm{sep}}(u)\) shows
\[
 \int_{\mathbb T^2}\{T(y)-y\}\,dy
 =\int_{\mathbb T^2}\{h_1^{\mathrm{sep}}(u)-h_{1\star}^{\mathrm{sep}}(u)\}
   \det Dh_{1\star}^{\mathrm{sep}}(u)\,du=0,                         \tag{7}
\]
because each component is a sine times products of constants and cosines.  Thus the selected translations of \(p\) and \(T\) are zero.  Their distance from the boundary of \(\mathcal Q\) is \(1/2\), and their graph distances from the cut locus are at least \(1/2-0.1/\pi\) and \(1/2-0.05/\pi\).  Choosing \(\beta\) below these margins proves lift regularity and therefore \(P^{\mathrm{sep}}\in\mathcal F_s\).

At \(u=(1/4,1/4)\), all cosines in (5) vanish.  With \(y=h_{1\star}^{\mathrm{sep}}(u)\),
\[
 y-T(y)=h_{1\star}^{\mathrm{sep}}(u)-h_1^{\mathrm{sep}}(u)
 =\left(\frac{0.05}{\pi},-\frac{0.05}{\pi}\right)\neq0.             \tag{8}
\]
At the same point \(H(y)=I_2\), \(f_{T1}(y)=1\), and \(\alpha(y)=1\).  Equal cell-size ratios give \(w_j=4\), and the covariance formula in \texttt{lem:full-critical-covariance} gives
\[
 V_j=\frac{4I_2}{4\pi}=\frac{I_2}{\pi}.                              \tag{9}
\]
This proves all claims of the proposition.
\end{proof}
\par\smallskip\noindent \textit{Justification.} The witness verifies nonemptiness, branch separation, nonzero displacement, and the candidate covariance at an analytic grid point. \textit{Gap.} Its separability does not replace the requested nonseparable moving-pole analysis. \textit{Consumer.} It supplies a reproducible benchmark for transport and adjoint solvers.

\begin{lemma}[lem:fixed-path-ma-linearization-source]\label{lem:fixed-path-ma-linearization-source}
For a fixed sufficiently regular source law and a spatially regular \(C^1\) path of positive target densities, the corresponding optimal-transport maps form a \(C^1\) path in a H\"older topology, and the path derivative is the uniquely normalized solution of the linearized Monge--Amp\`ere equation.
\end{lemma}

\begin{lemma}[lem:published-critical-boundary]\label{lem:published-critical-boundary}
\texttt{Proposition 29 of ManoleBalakrishnanNilesWeedWasserman2023CLT supplies the two-\allowbreak{}dimensional pointwise periodized-\allowbreak{}KDE optimal-\allowbreak{}transport Gaussian conclusion in the fixed-\allowbreak{}source uniform-\allowbreak{}to-\allowbreak{}uniform case;\allowbreak{} it does not establish the general variable-\allowbreak{}coefficient or generated-\allowbreak{}target case.}
\end{lemma}

\begin{lemma}[lem:uniform-critical-estimates]\label{lem:uniform-critical-estimates}
Under ass:density-positivity, ass:density-smoothness, ass:inner-ellipticity, ass:outer-ellipticity, ass:lift-regularity, ass:fixed-separated-grid, ass:kernel-symmetry, ass:kernel-mass, ass:kernel-order, ass:four-sample-design, and ass:polynomial-bandwidth, the following estimates hold uniformly over \(P\in\mathcal F_s\).  For \(h=h_n\), every coordinate of the score \(\Psi_{g,h}\) is the sum of its frozen-coefficient, first-order-coordinate Green score and a remainder whose \(L^2(P_g)\) norm is bounded by a common constant.  In particular,
\[
 \max_{g\in\mathcal G}\lVert\Psi_{g,h_n}\rVert_{L^2(P_g)}
 \leq C\sqrt{L_n},
 \qquad
 \max_{g\in\mathcal G}\mathbb E_{P_g}\lVert\Psi_{g,h_n}(X)\rVert^4
 \leq C h_n^{-2}.
\]
For the opposite-fold adjoint scores,
\[
 \max_{g\in\mathcal G,\ r\in\{1,2\}}
 L_n^{-1/2}
 \lVert\widehat\Psi_{g,h_n}^{(-r)}-\Psi_{g,h_n}\rVert_{L^2(P_g)}
 =o_{\mathbb P}(1),
\]
and, on events whose probabilities tend to one uniformly,
\[
 \max_{g,r}\mathbb E_{P_g}
 \lVert\widehat\Psi_{g,h_n}^{(-r)}(X)\rVert^4
 \leq C h_n^{-2}.
\]
Moreover, if \(u=(u_g)_{g\in\mathcal G}\) is a zero-integral four-density perturbation, the perturbed densities are positive, and \(\max_g\lVert u_g\rVert_{C^{2,1/4}}\) is smaller than a common constant, then the selected lifted nested map in def:identified-functional obeys
\[
 \max_{1\leq j\leq J}
 \left\lVert
 \widetilde T_{P+u}(\widetilde y_j)-\widetilde T_P(\widetilde y_j)
 -\mathcal D_P[u]_j
 \right\rVert
 \leq C\left(\max_{g\in\mathcal G}
 \lVert u_g\rVert_{C^{1,1/4}}\right)^2.
\]
For the signed periodized KDEs,
\[
 \sup_{P\in\mathcal F_s}
 \mathbb E_P\left[\max_g
 \lVert\widetilde f_g-f_g\rVert_{C^{1,1/4}}^4\right]
 \leq C\left\{
 \frac{\log n}{n h_n^{9/2}}+h_n^{2s-5/2}
 \right\}^2,
\]
and there are constants \(c_0,C_0,\epsilon_0>0\) such that
\[
 \sup_{P\in\mathcal F_s}\mathbb P_P\left(
 \max_g\lVert\widetilde f_g-f_g\rVert_{C^{2,1/4}}>\epsilon_0
 \right)
 \leq C_0\exp(-c_0 n h_n^{13/2}).
\]
On the complementary good event clipping is inactive, normalization is exactly one, the selected lift branches are the population branches, all population and pilot coefficients are uniformly elliptic, and every pilot pole differs from its population pole by \(o_{\mathbb P}(h_n)\).
\end{lemma}
\begin{proof}
Fix \(P\in\mathcal F_s\).  Constants below depend only on the displayed class constants, the fixed kernel, \(J\), and the fixed grid separation, and hence are uniform in \(P\).  Under the maximal bandwidth post-image, choose a fixed \(\vartheta\) satisfying
\[
 0<\vartheta<\min\left\{1,\frac{a_h^{-1}-6}{4}\right\}.
\tag{1}
\]
Under the narrower frozen bandwidth assumption the same choice is available.  Define
\[
 U_{k,n}=q_{k,n}(\vartheta)
 =\left\{\frac{\log n}{n h_n^{2+2k+2\vartheta}}\right\}^{1/2}
  +h_n^{s-k-\vartheta},\qquad k=0,1,2.
\tag{2}
\]
Writing the stochastic and bias summands in (2) separately gives
\[
 \sqrt{\frac n{L_n}}
 \left\{\frac{\log n}{n h_n^{4+2\vartheta}}\right\}^{1/2}
 \left\{\frac{\log n}{n h_n^{2+2\vartheta}}\right\}^{1/2}
 =a_h^{-1/2}\sqrt{\log n}\,
   n^{-1/2+a_h(3+2\vartheta)},
\]
while either stochastic--bias cross term equals
\(a_h^{-1/2}h_n^{s-2-2\vartheta}\), and the bias--bias term equals
\(n^{1/2-a_h(2s-1-2\vartheta)}/\sqrt{a_h\log n}\).
By (1), the first exponent is negative; by \(s\geq8\), the remaining powers vanish.  The same direct calculation gives
\(q_{2,n}(\vartheta)=o(1)\) and
\(q_{0,n}(\vartheta)/h_n=o(1)\).  Therefore
\[
 U_{2,n}=o(1),\qquad U_{0,n}/h_n=o(1),
 \qquad \sqrt{n/L_n}\,U_{1,n}U_{0,n}=o(1).
\tag{3}
\]

We first prove the density-pilot estimates.  For a multi-index \(b\) of length \(|b|=m\), a centered summand of
\(D^b\widetilde f_g(x)-\mathbb E_PD^b\widetilde f_g(x)\) has envelope at most
\(C h_n^{-2-m}\) and variance at most \(C h_n^{-2-2m}\); these follow directly by the changes of variables
\(u=(x-z)/h_n\), using the compact support and smoothness of \(K\), and ass:density-positivity.  Bernstein's exponential-moment argument on a torus lattice of mesh \(h_n n^{-3}\), followed by the mean-value theorem between lattice points, therefore yields, for \(m=0,1,2,3\),
\[
 \mathbb E_P\left\|D^m\widetilde f_g-
       \mathbb E_PD^m\widetilde f_g\right\|_\infty^4
 \leq C\left\{\frac{\log n}{n h_n^{2+2m}}\right\}^{2}.
\tag{4}
\]
The Bernstein linear term is absorbed because
\(n h_n^2/\log n\to\infty\), which follows from \(a_h<1/6\).  The deterministic interpolation inequality
\[
 [D^k v]_{C^{0,\vartheta}}
 \leq C\|D^k v\|_\infty^{1-\vartheta}
          \|D^{k+1}v\|_\infty^{\vartheta}
\tag{5}
\]
and H\"older's inequality applied to (4) give
\[
 \mathbb E_P\left\|\widetilde f_g-\mathbb E_P\widetilde f_g
                   \right\|_{C^{k,\vartheta}}^4
 \leq C\left\{\frac{\log n}
 {n h_n^{2+2k+2\vartheta}}\right\}^{2},
 \qquad k=0,1,2.
\tag{6}
\]
Taylor's formula with integral remainder, ass:kernel-mass, and every vanishing moment in ass:kernel-order give
\[
 \left\|\mathbb E_P\widetilde f_g-f_g\right\|_{C^{k,\vartheta}}
 \leq C h_n^{s-k-\vartheta},
 \qquad k=0,1,2.
\tag{7}
\]
Combining (6)--(7), taking the maximum over the four cells, and using
\((x+y)^4\leq8(x^4+y^4)\), we obtain
\[
 \mathbb E_P\left[\max_g
 \|\widetilde f_g-f_g\|_{C^{k,\vartheta}}^4\right]
 \leq C U_{k,n}^4.
\tag{8}
\]
The same lattice argument at a fixed threshold, now applied directly to H\"older increments, gives
\[
 \mathbb P_P\left(\max_g
 \|\widetilde f_g-f_g\|_{C^{2,\vartheta}}>\epsilon_0\right)
 \leq C_0\exp\{-c_0n h_n^{6+2\vartheta}\}
\tag{9}
\]
after decreasing \(\epsilon_0\) and taking \(n\) large enough that the bias in (7) is below \(\epsilon_0/2\).  Polynomial lattice factors are absorbed because (1) implies
\(n h_n^{6+2\vartheta}/\log n\to\infty\).  Setting \(\vartheta=1/4\) in (8) with \(k=1\) and in the unabsorbed form of (9) gives exactly
\[
 \mathbb E_P\left[\max_g
 \|\widetilde f_g-f_g\|_{C^{1,1/4}}^4\right]
 \leq C\left\{\frac{\log n}{n h_n^{9/2}}
                  +h_n^{2s-5/2}\right\}^{2},
\tag{10}
\]
and
\[
 \mathbb P_P\left(\max_g
 \|\widetilde f_g-f_g\|_{C^{2,1/4}}>\epsilon_0\right)
 \leq C_0\exp\{-c_0n h_n^{13/2}\},
\tag{11}
\]
as asserted in the target.  Notice that the high-probability event used below is (9) with the post-bandwidth choice (1); formula (11) remains the requested valid fixed-H\"older bound.

We next establish the periodic Green decomposition.  Fix a grid point \(y\), write \(B_y=B(y)\), and choose a coordinate ball whose radius is less than one eighth of the grid-separation and cut-locus margins.  In that chart set
\[
 \Gamma_y(z)=-\frac1{2\pi\sqrt{\det B_y}}
       \log\left|B_y^{-1/2}(z-y)\right|.
\tag{12}
\]
Let \(\chi\) be supported in the chart and equal to one near \(y\), and define before testing
\[
 F_y(z)=\{B(z)-B_y\}\chi(z)\nabla\Gamma_y(z).
\tag{13}
\]
Uniform H\"older control of the population coefficients, obtained from the uniformly invertible augmented systems in thm:nested-derivative, gives
\[
 |F_y(z)|\leq C|z-y|^{-1+\vartheta},
 \qquad \|F_y\|_{L^2}\leq C,
\tag{14}
\]
where the second inequality follows from
\(\int_0^r t^{-2+2\vartheta}t\,dt<\infty\).  The coercive functional on
\(H_0^1\) with linear part \(v\mapsto-\int F_y\cdot\nabla v\), together with the bounded cutoff terms, has a unique minimizer \(w_y\); testing its weak Euler equation is now legitimate and gives
\(\|\nabla w_y\|_2\leq C\).  The difference between the normalized periodic Green function and
\(\chi\Gamma_y+w_y\) is \(B\)-harmonic on a smaller ball.  The interior estimate in lem:elliptic-green-gradient-bound bounds its gradient there, while compactness and the mean-zero energy estimate bound it outside the chart.  Applying the same estimate after rescaling dyadic annuli in (13) yields
\[
 \nabla_zG_B(y,z)=\nabla\Gamma_y(z)+R_y(z),
 \qquad |R_y(z)|\leq C|z-y|^{-1+\vartheta},
 \qquad \sup_y\|R_y\|_2\leq C.
\tag{15}
\]
This construction supplies the torus harmonic correction rather than treating the bounded-domain Green estimate as periodic.

Insert a centered cell-
\(g\) kernel tangent into def:nested-derivative and apply Fubini before taking norms.  The coefficient identity
\(BD^{-\mathsf T}=\alpha\mathsf C\) proved in thm:nested-derivative makes the order-one part, in the original cell coordinate, equal to
\[
 \psi^0_{g,j,h}(x)
 =K_h^{\mathrm{per}}*\zeta^0_{g,j}(x)
   -P_g(K_h^{\mathrm{per}}*\zeta^0_{g,j}),
\tag{16}
\]
where \(\zeta^0_{g,j}\) is the affine pullback of \(\nabla\Gamma_{y_j}\) with the appropriate smooth cell weight.  The exact change of variables in the transformed integral cancels the Jacobian in each transformed density tangent; thus (16) retains both appearances of every control-cell noise.

For any population pullback \(\Phi\), Taylor's formula gives, uniformly for \(|v|\leq Ch\),
\[
 \Phi(x+v)-\Phi(x)-D\Phi(x)v
 =\int_0^1(1-t)D^2\Phi(x+tv)[v,v],dt,
 \qquad |\Phi(x+v)-\Phi(x)-D\Phi(x)v|\leq Ch^2.
\tag{17}
\]
Equations (15)--(17), split over dyadic annuli about the transformed pole, show that replacing the variable coefficient and nonlinear coordinate by the field in (16) costs at most a uniform constant in \(L^2(P_g)\).  Explicitly, away from a ball of radius \(Ch\), kernel symmetry and (17) give
\[
 |K_h^{\mathrm{per}}*\zeta^0(x)-\zeta^0(x)|
 \leq Ch,d(x,x_0)^{-2},
\tag{18}
\]
whereas inside the ball the mollified field is at most \(C/h\).  Consequently
\[
 \int_{d\leq Ch}|K_h^{\mathrm{per}}*\zeta^0|^2=O(1),
 \qquad
 \int_{d>Ch}|K_h^{\mathrm{per}}*\zeta^0-\zeta^0|^2=O(1).
\tag{19}
\]
The order-minus-two nested term is bounded by the convolution of two first-order Green singularities:
\[
 C\int_{\mathbb T^2}
 \frac{dz}{d(z,y_j)d(z,x)}
 \leq C\{1+|\log d(x,y_j)|\},
\tag{20}
\]
which is uniformly square integrable.  The weighted-Hodge construction and bounded Riesz scores in thm:nested-derivative give the same uniform \(L^2\) bound for the full harmonic part.  Since
\(\|K_h^{\mathrm{per}}\|_1\leq\|K\|_1^2\), Young's inequality is now applied to already established \(L^2\) remainders, and (15)--(20) yield
\[
 \|\Psi_{g,h}-\psi^0_{g,h}\|_{L^2(P_g)}\leq C,
 \qquad
 \|\Psi_{g,h}\|_{L^2(P_g)}\leq C\sqrt{\log(1/h)}.
\tag{21}
\]
Moreover, the annular envelope is \(C/(d(x,x_0)+h)\).  Because
\[
 \int_0^1\frac{r}{(r+h)^4},dr\leq Ch^{-2},
\tag{22}
\]
finite \(J\), density boundedness, and (20) imply
\[
 \max_g\mathbb E_{P_g}\|\Psi_{g,h}(X)\|^4\leq Ch^{-2}.
\tag{23}
\]
Centering is harmless because \(d(\cdot,x_0)^{-1}\) is integrable in two dimensions.  Equations (21)--(23) prove the first score assertions, including the generated-target cross terms and nonzero harmonic modes.

We next control estimated coefficients and poles.  On the opposite-fold event (9), the stability of the two augmented scalar--harmonic systems in thm:nested-derivative, the mixed one-map estimate in lem:one-map-mixed-linearization, and (8) imply
\[
 \epsilon_n^{\mathrm{coef}}
 :=\max_{g,r}\left(
 \|\widehat B^{(-r)}-B\|_{C^1}
 +\|\widehat{\mathsf C}^{(-r)}-\mathsf C\|_{C^1}
 +\|\widehat\alpha^{(-r)}-\alpha\|_{C^1}
 \right)=o_{\mathbb P}(1),
\tag{24}
\]
and every transformed evaluation pole satisfies
\[
 d_n^{\mathrm{pole}}=O_{\mathbb P}(U_{0,n})=o_{\mathbb P}(h_n).
\tag{25}
\]
Here the first equality in (25) follows by applying the uniformly bounded augmented inverse to the density perturbation in \(C^{0,\vartheta}\), and the last equality is (3).  For a mollified \(1/d\) field, differentiation in the pole and the two-region calculation in (19) give
\[
 \|K_h^{\mathrm{per}}*\zeta^0(\,\cdot-d)
       -K_h^{\mathrm{per}}*\zeta^0\|_2
 \leq C|d|/h.
\tag{26}
\]
For the scalar inverses the resolvent identity is
\[
 \mathcal L_{\widehat B}^{-1}-\mathcal L_B^{-1}
 =\mathcal L_{\widehat B}^{-1}
   (\mathcal L_B-\mathcal L_{\widehat B})\mathcal L_B^{-1};
\tag{27}
\]
the finite-dimensional Schur complements obey the identical algebraic identity.  Uniform coercivity from thm:nested-derivative, (21), and (24) show that coefficient changes cost at most
\(C\epsilon_n^{\mathrm{coef}}\sqrt{L_n}\) in score norm.  Equations (25)--(27), together with the bounded order-minus-two and harmonic envelopes, therefore give conditionally on the training fold and hence unconditionally
\[
 \max_{g,r}\,L_n^{-1/2}
 \|\widehat\Psi_{g,h_n}^{(-r)}-\Psi_{g,h_n}\|_{L^2(P_g)}
 =o_{\mathbb P}(1).
\tag{28}
\]
The first maximum in (28) is over \(g\in\mathcal G\) and the second over
\(r\in\{1,2\}\).  The estimated analog of the annular envelope in (22) holds on (9) with the same constant, so
\[
 \max_{g,r}\mathbb E_{P_g}
 \|\widehat\Psi_{g,h_n}^{(-r)}(X)\|^4\leq Ch_n^{-2}
\tag{29}
\]
on events with probability tending to one uniformly.

It remains to establish the nonlinear remainder with the correct two endpoints.  For a zero-integral perturbation \(u=(u_g)_g\), set
\[
 U_k(u)=\max_g\|u_g\|_{C^{k,\vartheta}},\qquad k=0,1,2.
\tag{30}
\]
Assume the perturbed densities are positive and \(U_2(u)<\epsilon_0\).  In the augmented equations from thm:nested-derivative, subtract the equation at \(P\) and its first derivative.  Reciprocal terms have the exact remainder
\[
 \frac1{f+u}-\frac1f+\frac{u}{f^2}
 =\frac{u^2}{f^2(f+u)},
\tag{31}
\]
bounded by \(C U_1(u)U_0(u)\) using positivity and
\(U_0(u)\leq U_1(u)\).  Product remainders contain one differentiated perturbation and one undifferentiated perturbation and have the same bound.  For compositions, the integral identity
\[
 v(x+d)-v(x)=\int_0^1Dv(x+td)d\,dt
\tag{32}
\]
shows that the post-linearization remainder is bounded by
\(C U_1(u)\|d\|_\infty+C\|d\|_\infty^2\); the augmented inverse gives
\(\|d\|_\infty\leq C U_0(u)\).  The determinant and scalar-potential terms are precisely the two-endpoint remainder controlled by lem:one-map-mixed-linearization, hence are also at most \(C U_1(u)U_0(u)\).  Finally, differentiating the weighted-Hodge reduced-cost stationarity equations shows that every harmonic residual is a bilinear integral with one coefficient or gradient perturbation and one displacement perturbation; uniform coercivity and Cauchy--Schwarz again bound it by \(C U_1(u)U_0(u)\).  Applying the uniformly bounded augmented inverse from thm:nested-derivative first to the inner system, propagating that solution through the exact generated density, and then applying it to the outer system gives
\[
 \max_{1\leq j\leq J}
 \left\|
 \widetilde T_{P+u}(\widetilde y_j)-\widetilde T_P(\widetilde y_j)
 -\mathcal D_P[u]_j
 \right\|
 \leq C U_1(u)U_0(u).
\tag{33}
\]
This is the required mixed \(C^{1,\vartheta}\)-times-
\(C^{0,\vartheta}\) estimate, including both harmonic translations.  It implies the frozen target's displayed squared \(C^{1,1/4}\) bound by taking \(\vartheta=1/4\) and using
\(U_0(u)\leq U_1(u)\).

For the KDE perturbation, (8), Cauchy--Schwarz, and (33) give the explicit probability and risk controls
\[
 \sqrt{n/L_n}\,U_{1,n}U_{0,n}=o(1),
 \qquad
 \mathbb E_P\|\operatorname{Rem}_n\|^2
 \leq C U_{1,n}^2U_{0,n}^2.
\tag{34}
\]
The deterministic order-\(s\) kernel bias is sent by the order-minus-two derivative to
\(O(h_n^{s-2})\), and (2) of the bandwidth proof makes
\(\sqrt{n/L_n}h_n^{s-2}\to0\).

Finally, on (9), decrease \(\epsilon_0\) below \(c/2\).  Then clipping is inactive.  Periodization and ass:kernel-mass give
\(\int_{\mathbb T^2}\widetilde f_g=1\), so the normalizer is exactly one.  Uniform augmented-system stability, ass:lift-regularity, and (24)--(25) keep both estimated maps on the selected population branches, preserve uniform ellipticity, and put every pilot pole at distance
\(o_{\mathbb P}(h_n)\) from its population pole.  Equations (8)--(11), (21)--(29), and (33)--(34) prove all assertions of the target, with the directed mixed remainder replacing the obsolete squared-rate calculation.
\end{proof}
\par\smallskip\noindent \textit{Justification.} The energy residual and weighted-Hodge harmonic block are now controlled separately. \textit{Gap.} Transformed kernels, estimated poles, fourth moments, opposite-fold stability, and the mixed two-endpoint remainder remain unresolved. \textit{Consumer.} This remains the analytic gate for covariance, Gaussian, ellipse, and upper-risk claims.

\begin{lemma}[lem:elliptic-green-gradient-bound]\label{lem:elliptic-green-gradient-bound}
Let \(A\) be a symmetric uniformly elliptic, Dini-continuous coefficient on a bounded planar domain.  Away from the boundary, the Green function of \(-\operatorname{div}(A\nabla\cdot)\) satisfies \(\lVert\nabla_xG_A(x,y)\rVert\leq C\lVert x-y\rVert^{-1}\), with \(C\) uniform when the ellipticity and Dini moduli are uniform.
\end{lemma}
\par\smallskip\noindent \textit{Justification.} The fourth-moment argument needs the pointwise pole bound separately from the energy parametrix. \textit{Gap.} This is classical elliptic substrate, not a nested transport contribution. \textit{Consumer.} Used in lem:uniform-critical-estimates.

\begin{lemma}[lem:one-map-mixed-linearization]\label{lem:one-map-mixed-linearization}
Fix \(b\in(0,1)\).  For positive densities \(p,q,\widehat q\in C^{2,b}(\mathbb T^2)\), with the source density \(p\) fixed, let \(\varphi_0\) and \(\widehat\varphi\) be the normalized Brenier potentials transporting \(p\) to \(q\) and \(\widehat q\), respectively, and let \(L\) be the linearized Monge--Amp\`ere operator at \((p,q)\).  Then
\[
 \left\lVert (\widehat\varphi-\varphi_0)-L^{-1}(\widehat q-q)\right\rVert_{C^{2,b}}
 \leq C\lVert\widehat q-q\rVert_{C^{1,b}}
          \lVert\widehat q-q\rVert_{C^{0,b}},
\]
where \(C\) is uniform on sets with uniform positive-density and H\"older-norm bounds.
\end{lemma}
\par\smallskip\noindent \textit{Justification.} This is the fixed-source mixed-norm template adapted explicitly to the two-endpoint nested construction. \textit{Gap.} The published theorem does not itself cover two varying endpoints or a generated target. \textit{Consumer.} Used in lem:uniform-critical-estimates.

\begin{lemma}[lem:mixed-bandwidth-feasibility]\label{lem:mixed-bandwidth-feasibility}
Under \(h_n=n^{-a_h}\), \(L_n=a_h\log n\), and \(s\geq8\), define
\[
 q_{2,n}=\sqrt{\frac{\log n}{n h_n^{13/2}}}+h_n^{s-9/4},\qquad
 q_{1,n}=\sqrt{\frac{\log n}{n h_n^{9/2}}}+h_n^{s-5/4},\qquad
 q_{0,n}=\sqrt{\frac{\log n}{n h_n^{5/2}}}+h_n^{s-1/4}.
\]
Then
\[
 \sqrt{n/L_n}\,h_n^{s-2}\longrightarrow0,\qquad
 q_{2,n}\longrightarrow0,\qquad
 \sqrt{n/L_n}\,q_{1,n}q_{0,n}\longrightarrow0
\]
hold simultaneously exactly when
\[
 \frac{1}{2(s-2)}\leq a_h<\frac17.
\]
In this window, \(n h_n^7/\log n\to\infty\), \(n h_n^{13/2}/\log n\to\infty\), and \(n h_n^2L_n^2\to\infty\).
\end{lemma}
\begin{proof}
For the parameterized calculation, fix \(0<\vartheta<1\) and put
\[
 q_{k,n}(\vartheta)
 =A_{k,n}(\vartheta)+D_{k,n}(\vartheta),\qquad
 A_{k,n}(\vartheta)
 =\left\{\frac{\log n}{n h_n^{2+2k+2\vartheta}}\right\}^{1/2},
 \qquad
 D_{k,n}(\vartheta)=h_n^{s-k-\vartheta},
 \quad k=0,1,2.
\tag{1}
\]
All summands in (1) are nonnegative.  Since \(h_n=n^{-a_h}\) and
\(L_n=a_h\log n\), the bias condition is
\[
 \sqrt{\frac n{L_n}}h_n^{s-2}
 =\frac{n^{1/2-a_h(s-2)}}{\sqrt{a_h\log n}}\longrightarrow0
 \quad\Longleftrightarrow\quad
 a_h\geq\frac1{2(s-2)}.
\tag{2}
\]
Equality is allowed in (2), because then the numerator's power of \(n\) is zero and the logarithmic denominator diverges.

The two requirements used for pilot and pole stability are weaker than the mixed requirement.  Indeed,
\[
 A_{2,n}(\vartheta)
 =\sqrt{\log n}\,n^{-1/2+a_h(3+\vartheta)},
 \qquad
 \frac{A_{0,n}(\vartheta)}{h_n}
 =\sqrt{\log n}\,n^{-1/2+a_h(2+\vartheta)}.
\tag{3}
\]
Thus \(q_{2,n}(\vartheta)=o(1)\) if
\(a_h<1/(6+2\vartheta)\), and
\(q_{0,n}(\vartheta)/h_n=o(1)\) if
\(a_h<1/(4+2\vartheta)\); their deterministic parts vanish because
\(s-2-\vartheta>0\) and \(s-1-\vartheta>0\) for \(s\geq8\) and
\(0<\vartheta<1\).

Expanding the product in the binding condition gives four nonnegative terms.  The stochastic--stochastic term is
\[
 \sqrt{\frac n{L_n}}A_{1,n}(\vartheta)A_{0,n}(\vartheta)
 =a_h^{-1/2}\sqrt{\log n}\,
   n^{-1/2+a_h(3+2\vartheta)}.
\tag{4}
\]
It converges to zero exactly when
\[
 a_h<\frac1{6+4\vartheta}.
\tag{5}
\]
The two stochastic--bias terms are identical after normalization:
\[
 \sqrt{\frac n{L_n}}A_{1,n}(\vartheta)D_{0,n}(\vartheta)
 =\sqrt{\frac n{L_n}}D_{1,n}(\vartheta)A_{0,n}(\vartheta)
 =a_h^{-1/2}h_n^{s-2-2\vartheta}=o(1).
\tag{6}
\]
The last equality follows from \(s\geq8\) and \(\vartheta<1\).  Finally,
\[
 \sqrt{\frac n{L_n}}D_{1,n}(\vartheta)D_{0,n}(\vartheta)
 =\frac{n^{1/2-a_h(2s-1-2\vartheta)}}{\sqrt{a_h\log n}}.
\tag{7}
\]
Under the lower bound in (2), the power in (7) is strictly negative, since
\(2s-1-2\vartheta>s-2\).  Hence (5) is the only upper restriction.  It also implies both upper restrictions following (3), because
\(6+4\vartheta>6+2\vartheta\) and
\(6+4\vartheta>4+2\vartheta\).

Now fix \(a_h\).  If
\[
 \frac1{2(s-2)}\leq a_h<\frac16,
\tag{8}
\]
then \(a_h^{-1}-6>0\), so one may choose once and for all
\[
 0<\vartheta<\min\left\{1,\frac{a_h^{-1}-6}{4}\right\}.
\tag{9}
\]
For this choice, (5) holds, and (2)--(7) prove
\[
 \sqrt{n/L_n}\,h_n^{s-2}\to0,
 \qquad q_{2,n}(\vartheta)\to0,
 \qquad q_{0,n}(\vartheta)/h_n\to0,
 \qquad
 \sqrt{n/L_n}\,q_{1,n}(\vartheta)q_{0,n}(\vartheta)\to0.
\tag{10}
\]
Conversely, if (10) holds for some fixed \(\vartheta>0\), (2) forces the lower inequality in (8), while the nonnegative term (4) forces
\(a_h<1/(6+4\vartheta)<1/6\).  At \(a_h=1/6\), (4) equals
\(a_h^{-1/2}\sqrt{\log n}\,n^{\vartheta/3}\) and diverges.  This proves that (8), with the upper endpoint excluded, is the maximal polynomial window when \(\vartheta\) is selected after \(a_h\).

The statement in the frozen core is the specialization \(\vartheta=1/4\).  In that case
\[
 q_{2,n}=q_{2,n}(1/4),\qquad
 q_{1,n}=q_{1,n}(1/4),\qquad
 q_{0,n}=q_{0,n}(1/4),
\]
and (5) becomes \(a_h<1/7\).  Equations (2)--(7) therefore prove exactly the three equivalences asserted there.  In this specialized window,
\[
 \frac{n h_n^7}{\log n}
 =\frac{n^{1-7a_h}}{\log n}\to\infty,
 \qquad
 \frac{n h_n^{13/2}}{\log n}
 =\frac{n^{1-(13/2)a_h}}{\log n}\to\infty,
\tag{11}
\]
because \(a_h<1/7<2/13\), and
\[
 n h_n^2L_n^2
 =a_h^2n^{1-2a_h}(\log n)^2\to\infty
\tag{12}
\]
because \(a_h<1/2\).  This proves every assertion of the target and, in addition, the directed maximal \(1/6\) parameterized window.
\end{proof}
\par\smallskip\noindent \textit{Justification.} Direct exponent arithmetic proves the mixed high-times-low feasibility window. \textit{Gap.} The broader endpoint is arithmetic only and is not adopted by the estimator assumptions until the coupled mixed remainder itself is proved. \textit{Consumer.} It records the scale available to a future repaired uniform expansion without overstating current inference.

\begin{theorem}[thm:uniform-critical-control]\label{thm:uniform-critical-control}
The answer to the open question is affirmative on the maintained class and on the full open bandwidth window.  Fix
\[
 \frac{1}{2(s-2)}\leq a_h<\frac16
\]
and choose once and for all
\[
 0<\vartheta<\min\left\{1,\frac{a_h^{-1}-6}{4}\right\}.
\]
Put
\[
 q_{k,n}(\vartheta)=
 \left\{\frac{\log n}{n h_n^{2+2k+2\vartheta}}\right\}^{1/2}
 +h_n^{s-k-\vartheta},\qquad k=0,1,2.
\]
There is a finite constant \(C_{\mathrm{uc}}\), depending only on the maintained class constants, the kernel, and the fixed grid, for which the following assertions hold.
Uniformly over \(P\in\mathcal F_s\), every full nested score admits a decomposition
\[
 \Psi_{g,h_n}=\psi^0_{g,h_n}+R_{g,h_n},
 \qquad
 \max_{g\in\mathcal G}\lVert R_{g,h_n}\rVert_{L^2(P_g)}\leq C_{\mathrm{uc}},
\]
where \(\psi^0_{g,h_n}\) is the centered score obtained by freezing the outer coefficient at each fixed Green pole and replacing every smooth pullback by its first-order affine coordinate map.  Consequently, if
\[
 \Sigma^0_{h_n}(P)=L_n^{-1}
 \sum_{g\in\mathcal G}\rho_g^{-1}
 \operatorname{Cov}_{P_g}(\psi^0_{g,h_n}(X)),
\]
then
\[
 \sup_{P\in\mathcal F_s}
 \lVert\Sigma_{h_n}(P)-\Sigma^0_{h_n}(P)\rVert_{\mathrm{op}}
 \longrightarrow0.
\]
The score generated by the second, two-elliptic-inverse term of \(\mathcal D_P\) and the score generated by the complete scalar--harmonic correction \(\mathcal H_P\), including both translation modes, have uniformly bounded \(L^2(P_g)\) norms.  For the opposite-fold feasible scores,
\[
 \sup_{P\in\mathcal F_s}\mathbb P_P\!\left\{
 \max_{g\in\mathcal G,\ r\in\{1,2\}}
 L_n^{-1/2}
 \lVert\widehat\Psi_{g,h_n}^{(-r)}-\Psi_{g,h_n}\rVert_{L^2(P_g)}
 >\varepsilon\right\}\longrightarrow0
\]
for every \(\varepsilon>0\).  On events whose probabilities tend to one uniformly, the opposite-fold estimated coefficients are uniformly elliptic and converge to their population coefficients, every transformed estimated pole is \(O_{\mathbb P}(q_{0,n}(\vartheta))=o_{\mathbb P}(h_n)\) from its population pole, and the selected inner and outer lift branches equal their population branches.

Finally, for every zero-integral four-density perturbation \(u=(u_g)_{g\in\mathcal G}\) with positive perturbed densities, define
\[
 U_k(u)=\max_{g\in\mathcal G}\lVert u_g\rVert_{C^{k,\vartheta}},
 \qquad k=0,1,2.
\]
There is an \(\epsilon_0>0\), uniform over \(P\in\mathcal F_s\), such that \(U_2(u)\leq\epsilon_0\) implies the translation-aware mixed remainder bound
\[
 \max_{1\leq j\leq J}
 \left\lVert
 \widetilde T_{P+u}(\widetilde y_j)-\widetilde T_P(\widetilde y_j)
 -\mathcal D_P[u]_j
 \right\rVert
 \leq C_{\mathrm{uc}} U_1(u)U_0(u).
\]
For the periodized-KDE perturbation, \(U_k=O_{\mathbb P}(q_{k,n}(\vartheta))\) uniformly.  Denote the vector on the left side of the preceding display, evaluated at \(u_g=\widetilde f_g-f_g\), by \(\operatorname{Rem}_n\).  Then
\[
 \sqrt{\frac n{L_n}}q_{1,n}(\vartheta)q_{0,n}(\vartheta)\to0,
 \qquad
 \sqrt{\frac n{L_n}}\operatorname{Rem}_n=o_{\mathbb P}(1),
 \qquad
 \sup_{P\in\mathcal F_s}\mathbb E_P
 \lVert\operatorname{Rem}_n\rVert^2
 \leq C_{\mathrm{uc}}q_{1,n}^2(\vartheta)q_{0,n}^2(\vartheta).
\]
Thus the transformed-kernel covariance replacement, the bounded second-elliptic and complete harmonic scores, the coefficient, moving-pole, and selected-lift stability, and the nonlinear critical negligibility requested in \(\mathfrak H_{\mathrm{crit}}\) all hold; the endpoint \(a_h=1/6\) is excluded.
\end{theorem}
\begin{proof}
Fix \(P\in\mathcal F_s\).  By the definition of \(\mathcal F_s\), the four densities are bounded below by \(c>0\), have uniformly bounded \(C^s\) norms, the inner and outer Jacobians are uniformly elliptic, and the selected lift branches have a positive margin.  Constants below therefore depend only on the maintained class constants, the fixed kernel, the fixed number \(J\), and the fixed grid separation.  In particular, they do not depend on \(P\).

First consider the bandwidth arithmetic.  Since \(a_h<1/6\), one has \(a_h^{-1}-6>0\), so the displayed interval for \(\vartheta\) is nonempty.  Write
\[
 A_{k,n}=\left\{\frac{\log n}
 {n h_n^{2+2k+2\vartheta}}\right\}^{1/2},
 \qquad
 D_{k,n}=h_n^{s-k-\vartheta},
 \qquad q_{k,n}=A_{k,n}+D_{k,n}.
\tag{1}
\]
The stochastic--stochastic part of the normalized mixed product is
\[
 \sqrt{\frac n{L_n}}A_{1,n}A_{0,n}
 =a_h^{-1/2}\sqrt{\log n}\,
 n^{-1/2+a_h(3+2\vartheta)}=o(1),
\tag{2}
\]
because the choice of \(\vartheta\) is equivalent to
\(a_h(6+4\vartheta)<1\).  The two stochastic--bias terms satisfy
\[
 \sqrt{\frac n{L_n}}A_{1,n}D_{0,n}
 =\sqrt{\frac n{L_n}}D_{1,n}A_{0,n}
 =a_h^{-1/2}h_n^{s-2-2\vartheta}=o(1),
\tag{3}
\]
since \(s\geq8\) and \(\vartheta<1\).  The bias--bias term is
\[
 \sqrt{\frac n{L_n}}D_{1,n}D_{0,n}
 =\frac{n^{1/2-a_h(2s-1-2\vartheta)}}{\sqrt{a_h\log n}}=o(1).
\tag{4}
\]
Indeed, \(a_h\geq1/[2(s-2)]\) and
\(2s-1-2\vartheta>s-2\), so the power of \(n\) in (4) is strictly negative.  The same inequalities give
\[
 q_{2,n}=o(1),\qquad q_{0,n}/h_n=o(1),
 \qquad \sqrt{n/L_n}\,q_{1,n}q_{0,n}=o(1).
\tag{5}
\]
These are precisely the parameterized conclusions proved in
\texttt{lem:mixed-bandwidth-feasibility}.  If \(a_h=1/6\), the expression in (2) is
\(a_h^{-1/2}\sqrt{\log n}\,n^{\vartheta/3}\), which diverges for every \(\vartheta>0\); this proves the asserted exclusion of the endpoint.

We next record the uniform pilot event used below.  A derivative of order \(k\) of a two-dimensional periodized kernel summand has envelope \(Ch_n^{-2-k}\) and variance at most \(Ch_n^{-2-2k}\).  Bernstein's inequality on an \(h_n n^{-3}\)-mesh of the torus, the mean-value theorem between mesh points, and the interpolation inequality
\[
 [D^k v]_{C^{0,\vartheta}}
 \leq C\lVert D^k v\rVert_\infty^{1-\vartheta}
          \lVert D^{k+1}v\rVert_\infty^\vartheta
\tag{6}
\]
give, for \(k=0,1,2\),
\[
 \sup_{P\in\mathcal F_s}\mathbb E_P\!\left[
 \max_{g\in\mathcal G}
 \lVert\widetilde f_g-f_g\rVert_{C^{k,\vartheta}}^4\right]
 \leq Cq_{k,n}^4.
\tag{7}
\]
Here the deterministic term in (7) follows from Taylor's formula, the unit kernel mass, the vanishing moments through order \(s-1\), and the uniform \(C^s\) bound.  The same mesh argument at a fixed threshold yields
\[
 \sup_{P\in\mathcal F_s}\mathbb P_P\!\left\{
 \max_g\lVert\widetilde f_g-f_g\rVert_{C^{2,\vartheta}}>epsilon_0
 \right\}
 \leq C_0\exp\{-c_0nh_n^{6+2\vartheta}\}=o(1),
\tag{8}
\]
because \(a_h(6+2\vartheta)<a_h(6+4\vartheta)<1\).  Equations (6)--(8) are also the parameterized pilot bounds in
\texttt{lem:uniform-critical-estimates}.

We now prove the frozen-pole covariance replacement.  By
\texttt{lem:uniform-critical-estimates}, application of the exact derivative from
\texttt{thm:nested-derivative} to a centered cell kernel tangent gives, after the exact pullback changes of variables,
\[
 \Psi_{g,h_n}=\psi^0_{g,h_n}+R_{g,h_n},
 \qquad
 \lVert R_{g,h_n}\rVert_{L^2(P_g)}\leq C,
 \qquad
 \lVert\Psi_{g,h_n}\rVert_{L^2(P_g)}\leq C\sqrt{L_n}.
\tag{9}
\]
The first term in (9) uses the anisotropic fundamental solution at a grid point \(y_j\),
\[
 \Gamma_j(z)=-\frac{1}{2\pi\sqrt{\det B(y_j)}}
 \log\left|B(y_j)^{-1/2}(z-y_j)\right|,
\tag{10}
\]
and the affine linearization of the relevant cell pullback.  Positivity and ellipticity make (10) well defined.  The torus Green-parametrix argument in
\texttt{lem:uniform-critical-estimates}, based on the interior estimate in
\texttt{lem:elliptic-green-gradient-bound}, gives the bounded remainder in (9); the fixed separation of the grid keeps distinct poles uniformly apart.

By the triangle inequality, (9) also gives
\(\lVert\psi^0_{g,h_n}\rVert_{L^2(P_g)}\leq C\sqrt{L_n}\).  For centered vector scores, Cauchy--Schwarz and (9) imply
\[
 \begin{aligned}
 &\left\lVert
 \operatorname{Cov}_{P_g}(\Psi_{g,h_n}(X))
 -\operatorname{Cov}_{P_g}(\psi^0_{g,h_n}(X))
 \right\rVert_{\mathrm{op}}\\
 &\qquad\leq
 2\lVert\psi^0_{g,h_n}\rVert_{L^2(P_g)}
   \lVert R_{g,h_n}\rVert_{L^2(P_g)}
 +\lVert R_{g,h_n}\rVert_{L^2(P_g)}^2
 \leq C\sqrt{L_n}.
\end{aligned}
\tag{11}
\]
There are four cells, their limiting ratios lie in \([1,\bar\rho]\), and \(L_n\to\infty\).  Dividing (11) by \(L_n\), multiplying by \(\rho_g^{-1}\), and summing proves
\[
 \sup_{P\in\mathcal F_s}
 \lVert\Sigma_{h_n}(P)-\Sigma^0_{h_n}(P)\rVert_{\mathrm{op}}
 \leq C L_n^{-1/2}\longrightarrow0.
\tag{12}
\]
Thus no pointwise-consistency argument is being used to move a singular pole: (12) follows from the uniform score comparison (9).

For completeness, isolate the two bounded-score components.  The kernel score of the second elliptic term in \(\mathcal D_P\) is bounded, up to smooth uniformly bounded weights, by a convolution of two Green gradients.  The Green-gradient bound gives
\[
 |\zeta_{2,j}(x)|
 \leq C\int_{\mathbb T^2}
 \frac{dz}{d_{\mathbb T^2}(y_j,z)d_{\mathbb T^2}(z,x)}
 \leq C\{1+|\log d_{\mathbb T^2}(x,y_j)|\}.
\tag{13}
\]
To verify the last inequality, set \(r=d_{\mathbb T^2}(x,y_j)\), split into the radius-\(r/2\) balls about \(x\) and \(y_j\), and use
\((ab)^{-1}\leq(a^{-2}+b^{-2})/2\) on the complement.  Polar integration then gives a constant on the two balls and \(C\log(1/r)\) on the complement.  Since
\[
 \int_0^{1/2}\{1+|\log r|\}^2r\,dr<\infty,
\tag{14}
\]
the right-hand side of (13) is uniformly square integrable.  Convolution with the periodized signed kernel preserves this bound because
\(\lVert K_{h_n}^{\mathrm{per}}\rVert_1\leq\lVert K\rVert_1^2\).  Hence the second elliptic score is uniformly bounded in \(L^2(P_g)\).  The complete harmonic score is uniformly bounded in \(L^2(P_g)\) by the Riesz-score conclusion of
\texttt{thm:nested-derivative}; that conclusion includes the inner and outer translation increments and both coefficient-divergence correctors.  This proves the bounded-score assertions without deleting any generated-target or harmonic term.

We next address feasible coefficients and moving poles.  On the event in (8), stability of the two uniformly invertible augmented scalar--harmonic systems in
\texttt{thm:nested-derivative} and the coefficient estimates in
\texttt{lem:uniform-critical-estimates} give convergence of every opposite-fold coefficient in \(C^1\).  The corresponding pole displacement is
\[
 d_n^{\mathrm{pole}}=O_{\mathbb P}(q_{0,n})=o_{\mathbb P}(h_n)
\tag{15}
\]
by (5).  Translation of a kernel-mollified \(1/d\) field by a vector \(d\) costs at most \(C|d|/h_n\) in \(L^2\).  Writing \(\widehat B\) for a generic opposite-fold estimate of the outer coefficient, the elliptic resolvent identity is
\[
 \mathcal L_{\widehat B}^{-1}-\mathcal L_B^{-1}
 =\mathcal L_{\widehat B}^{-1}
   (\mathcal L_B-\mathcal L_{\widehat B})\mathcal L_B^{-1}
\tag{16}
\]
and the identical finite-dimensional Schur-complement identity show that a coefficient error \(\epsilon_n=o_{\mathbb P}(1)\) costs at most
\(C\epsilon_n\sqrt{L_n}\) in score norm.  Combining (9), (15), and (16), exactly as in
\texttt{lem:uniform-critical-estimates}, yields for every \(\varepsilon>0\)
\[
 \sup_{P\in\mathcal F_s}\mathbb P_P\!\left\{
 \max_{g,r}L_n^{-1/2}
 \lVert\widehat\Psi_{g,h_n}^{(-r)}-\Psi_{g,h_n}\rVert_{L^2(P_g)}
 >\varepsilon\right\}\longrightarrow0.
\tag{17}
\]
The positive lift margin and (8), (15) keep both estimated maps in the same smooth cost chart and on the population-selected branches.  The complement has probability tending to zero uniformly.  This proves all coefficient, pole, and selected-lift claims in the statement.

It remains to verify the mixed nonlinear remainder.  Let \(u\) satisfy the conditions in the statement.  Subtract from each inner or outer augmented mass-and-stationarity system its value at \(P\) and its first derivative.  Positivity gives the exact reciprocal remainder
\[
 \frac1{f+u}-\frac1f+\frac{u}{f^2}
 =\frac{u^2}{f^2(f+u)},
\tag{18}
\]
whose relevant H\"older norm is at most \(CU_1(u)U_0(u)\).  Every product remainder contains one differentiated and one undifferentiated perturbation and has the same bound.  For compositions, the identity
\[
 v(x+d)-v(x)=\int_0^1Dv(x+td)d\,dt
\tag{19}
\]
and the uniformly bounded augmented inverse give
\(\lVert d\rVert_\infty\leq CU_0(u)\); the post-linearization part of (19) is therefore bounded by
\[
 CU_1(u)\lVert d\rVert_\infty+C\lVert d\rVert_\infty^2
 \leq CU_1(u)U_0(u).
\tag{20}
\]
The determinant and scalar-potential residuals obey the same two-endpoint bound by the mixed one-map estimate in
\texttt{lem:one-map-mixed-linearization}, applied after each exact endpoint pullback.  The change of the pullback itself is controlled by (19), so this application does not treat the generated target as fixed.  Finally, differentiating the weighted-Hodge stationarity equations leaves bilinear residuals with one coefficient or gradient perturbation and one displacement perturbation; their norms are bounded by \(CU_1(u)U_0(u)\) by Cauchy--Schwarz and uniform coercivity from
\texttt{thm:nested-derivative}.

Apply the uniformly bounded augmented inverse first to the inner residual, propagate its solution through the exact generated density using (18)--(20), and then apply the outer augmented inverse.  This gives
\[
 \max_{1\leq j\leq J}
 \left\lVert
 \widetilde T_{P+u}(\widetilde y_j)-\widetilde T_P(\widetilde y_j)
 -\mathcal D_P[u]_j
 \right\rVert
 \leq CU_1(u)U_0(u).
\tag{21}
\]
Both harmonic translations are included because the augmented inverses used in this argument are the scalar--harmonic inverses of
\texttt{thm:nested-derivative}.

For \(u_g=\widetilde f_g-f_g\), equations (7) and (21), followed by Cauchy--Schwarz, yield
\[
 \sup_{P\in\mathcal F_s}\mathbb E_P
 \lVert\operatorname{Rem}_n\rVert^2
 \leq C
 \left(\mathbb E_PU_1^4\right)^{1/2}
 \left(\mathbb E_PU_0^4\right)^{1/2}
 \leq Cq_{1,n}^2q_{0,n}^2.
\tag{22}
\]
This calculation is first made on the event (8).  On its complement, compactness of the torus bounds the selected displacement and its fixed fallback, while the deterministic kernel derivative envelopes bound the linear score and all algebraic plug-in terms by a fixed polynomial in \(h_n^{-1}\).  Because the probability in (8) is \(\exp\{-c n^{1-a_h(6+2\vartheta)}\}\) with a strictly positive exponent, Cauchy--Schwarz shows that the complement contributes less than every fixed polynomial rate.  It is therefore absorbed by the right-hand side of (22).
Equations (5), (21), and Markov's inequality imply
\[
 \sqrt{n/L_n}\,\operatorname{Rem}_n=o_{\mathbb P}(1)
\tag{23}
\]
uniformly over \(\mathcal F_s\).  The deterministic kernel bias is mapped by the order-minus-two derivative to \(O(h_n^{s-2})\), and
\[
 \sqrt{n/L_n}\,h_n^{s-2}
 =\frac{n^{1/2-a_h(s-2)}}{\sqrt{a_h\log n}}\longrightarrow0
\tag{24}
\]
by the lower bandwidth bound, including equality.  On (8), decreasing \(\epsilon_0\) below \(c/2\) makes clipping inactive; unit kernel mass makes the normalizer exactly one.  Hence (21)--(24) apply to the specified positive-normalized plug-in, while the fallback event has uniformly vanishing probability.  Equations (12)--(17) and (21)--(24) establish every clause of the affirmative answer.
\end{proof}
\par\smallskip\noindent \textit{Justification.} The parameterized Green-parametrix, opposite-fold stability, and mixed scalar--harmonic expansion close every clause of the former analytic open question on the widened bandwidth window. \textit{Gap.} There is no remaining internal analytic gap under the maintained class; the excluded endpoint and broader data regimes remain outside scope. \textit{Consumer.} This theorem supplies the singular-score and nonlinear-remainder gate used by covariance, Gaussian, adjoint, ellipse, and risk results.

\section{Interpretation}

At each prespecified treated-post location, the estimand is the selected-lift displacement under the identified no-treatment map.  The logarithmic variance combines all four cells, including the generated-target cross products within each control cell; translation and second-elliptic uncertainty are retained at finite bandwidth but vanish after logarithmic normalization.  The proved Gaussian and adjoint results validate the simultaneous ellipses.  Because the limiting grid blocks are independent and each standardized quadratic form is \(\chi^2_2\), the deterministic limiting critical value is \(c_{J,\gamma}=-2\log\{1-(1-\gamma)^{1/J}\}\).

\section*{Technical internal limitation (diagnostic only)}

The result is uniform only for the maintained order-\(s\), \(s\geq8\), positive-density, uniformly elliptic, selected-no-cut flat-torus class, fixed separated \(J\), independent four-cell samples, and polynomial bandwidth \(1/[2(s-2)]\leq a_h<1/6\).  It does not cover the endpoint \(a_h=1/6\), cut-locus crossings, dependent panels, growing grids, nonperiodic domains, or unsmoothed empirical optimal transport, and it does not establish semiparametric efficiency or an exact minimax constant.

\section{Honest scope}

Within the stated torus class, the covariance, central limit theorem, adjoint covariance estimator, simultaneous multiplier ellipses, and \(L_n/n\) upper and Hellinger-local lower squared-risk rates are proved.  The lower bound matches the rate but not an exact constant.  A Card--Krueger bivariate employment reanalysis still requires a prespecified defensible periodic outcome embedding, assessment of the independent-cell repeated-cross-section approximation, and numerical transport and adjoint errors below the \(\sqrt{L_n/n}\) statistical scale; none of those empirical design choices is supplied by the theorem.

\begin{thebibliography}{99}

  \bibitem{TorousGunsiliusRigollet2024OTCiC} William Torous, Florian Gunsilius, and Philippe Rigollet (2024), ``An Optimal Transport Approach to Estimating Causal Effects via Nonlinear Difference-in-Differences,'' Journal of Causal Inference 12; arXiv:2108.05858.

  \bibitem{AtheyImbens2006CiC} Susan Athey and Guido W. Imbens (2006), ``Identification and Inference in Nonlinear Difference-in-Differences Models,'' Econometrica 74, 431--497.

  \bibitem{ManoleBalakrishnanNilesWeedWasserman2023CLT} Tudor Manole, Sivaraman Balakrishnan, Jonathan Niles-Weed, and Larry Wasserman (2023), ``Central Limit Theorems for Smooth Optimal Transport Maps,'' arXiv:2312.12407.

  \bibitem{GonzalezSanzSheng2024Linearization} Alberto Gonzalez-Sanz and Shunan Sheng (2024), ``Linearization of Monge--Ampere Equations and Statistical Applications,'' arXiv:2408.06534.

  \bibitem{GonzalezSanzShengWuAvellaMedina2026Influence} Alberto Gonzalez-Sanz, Shunan Sheng, Bohan Wu, and Marco Avella Medina (2026), ``The Influence Function of Transport-Based Quantiles,'' arXiv:2607.19080v1.

  \bibitem{GonzalezSanzMordantSheng2026PotentialFCLT} Alberto Gonzalez-Sanz, Gilles Mordant, and Shunan Sheng (2026), ``Empirical Optimal Transport Potentials: Fast Rates and a Functional Central Limit Theorem,'' arXiv:2608.00649.

  \bibitem{HutterRigollet2021Minimax} Jan-Christian Hutter and Philippe Rigollet (2021), ``Minimax Estimation of Smooth Optimal Transport Maps,'' Annals of Statistics 49, 1166--1194.

  \bibitem{ManoleBalakrishnanNilesWeedWasserman2024Plugin} Tudor Manole, Sivaraman Balakrishnan, Jonathan Niles-Weed, and Larry Wasserman (2024), ``Plugin Estimation of Smooth Optimal Transport Maps,'' Annals of Statistics 52, 966--998.

  \bibitem{PonnopratOkanoImaizumi2024UniformBand} Donlapark Ponnoprat, Ryo Okano, and Masaaki Imaizumi (2024), ``Uniform Confidence Band for Optimal Transport Map on One-Dimensional Data,'' Electronic Journal of Statistics 18, 515--546.

  \bibitem{SadhuGoldfeldKato2023SemidiscreteInference} Ritwik Sadhu, Ziv Goldfeld, and Kengo Kato (2023), ``Stability and Statistical Inference for Semidiscrete Optimal Transport Maps,'' arXiv:2303.10155.

  \bibitem{AuricchioVeneroni2022DiscreteOT} Gennaro Auricchio and Marco Veneroni (2022), ``On the Structure of Optimal Transportation Plans Between Discrete Measures,'' Applied Mathematics and Optimization 85, article 42.

  \bibitem{FangSantos2019DirectionalInference} Zheng Fang and Andres Santos (2019), ``Inference on Directionally Differentiable Functions,'' Review of Economic Studies 86, 377--412.

  \bibitem{ChernozhukovChetverikovKato2013Multiplier} Victor Chernozhukov, Denis Chetverikov, and Kengo Kato (2013), ``Gaussian Approximations and Multiplier Bootstrap for Maxima of Sums of High-Dimensional Random Vectors,'' Annals of Statistics 41, 2786--2819.

  \bibitem{SunTchetgenTchetgen2025DebiasedCiC} Jinghao Sun and Eric J. Tchetgen Tchetgen (2025), ``On a Debiased and Semiparametric Efficient Changes-in-Changes Estimator,'' arXiv:2507.07228v2.

  \bibitem{McCann2001PolarFactorization} Robert J. McCann (2001), ``Polar Factorization of Maps on Riemannian Manifolds,'' Geometric and Functional Analysis 11, 589--608.

  \bibitem{GruterWidman1982GreenFunction} Michael Gruter and Kjell-Ove Widman (1982), ``The Green Function for Uniformly Elliptic Equations,'' Manuscripta Mathematica 37, 303--342, doi:10.1007/BF01166225.

\end{thebibliography}

\end{document}